% Capítulo 52 del sucesor: integración pública del propietario íntegro del radión.
% No carga portada, main, style ni C_estatutos_procedencia del módulo autónomo.

% Adaptador de compatibilidad del propietario íntegro del radión.
%
% El tratado rector ya carga tipografía, geometría, hipervínculos, teoremas,
% tablas, TikZ y tcolorbox.  Este archivo aporta únicamente los nombres y
% entornos exclusivos del propietario, sin importar su preámbulo autónomo ni
% alterar la jerarquía tipográfica general.

\makeatletter
\@ifundefined{color@RadNavy}{\definecolor{RadNavy}{HTML}{102A43}}{}
\@ifundefined{color@RadBlue}{\definecolor{RadBlue}{HTML}{1F5E7A}}{}
\@ifundefined{color@RadTeal}{\definecolor{RadTeal}{HTML}{168C8C}}{}
\@ifundefined{color@RadCopper}{\definecolor{RadCopper}{HTML}{B66A3C}}{}
\@ifundefined{color@RadGold}{\definecolor{RadGold}{HTML}{D4A73A}}{}
\@ifundefined{color@RadInk}{\definecolor{RadInk}{HTML}{1A202C}}{}
\@ifundefined{color@RadMist}{\definecolor{RadMist}{HTML}{EDF4F7}}{}
\@ifundefined{color@RadCream}{\definecolor{RadCream}{HTML}{FAF7F0}}{}
\@ifundefined{color@RadRose}{\definecolor{RadRose}{HTML}{8D3F55}}{}
\@ifundefined{color@RadGreen}{\definecolor{RadGreen}{HTML}{2F7D5D}}{}
\@ifundefined{color@RadGray}{\definecolor{RadGray}{HTML}{66788A}}{}

\providecommand{\TRIT}{\ensuremath{\mathrm{TRIT}}}
\providecommand{\HMT}{\ensuremath{\mathrm{HMT}}}
\providecommand{\Rstate}{\mathcal R_{12}}
\providecommand{\epsR}{\varepsilon_R}
\providecommand{\epsone}{\varepsilon_R^{(1)}}
\providecommand{\pih}{\pi_{\!\HMT}}
\providecommand{\alphah}{\alpha_{\!\HMT}}
\providecommand{\hret}{\hbar_{\mathrm{ret}}}
\providecommand{\hfull}{h_{\mathrm{ret}}}
\providecommand{\diff}{\mathop{}\!\mathrm d}
\providecommand{\e}{\mathrm e}
\providecommand{\R}{\mathbb R}
\providecommand{\C}{\mathbb C}
\providecommand{\F}{\mathbb F}
\providecommand{\Z}{\mathbb Z}
\providecommand{\dd}{\,\mathrm d}
\providecommand{\status}[1]{\textsf{\bfseries #1}}
\providecommand{\sourcepath}[1]{\nolinkurl{#1}}

\@ifundefined{typebox}{%
  \newtcolorbox{typebox}[1][]{enhanced,breakable,colback=white,
    colframe=RadBlue,boxrule=.7pt,arc=1.5mm,left=4mm,right=4mm,
    top=3mm,bottom=3mm,title={Recibo de tipos},fonttitle=\bfseries,
    coltitle=RadNavy,#1}%
}{}
\@ifundefined{narrativebox}{%
  \newtcolorbox{narrativebox}[1][]{enhanced,breakable,colback=white,
    colframe=RadGold,borderline west={2.4pt}{0pt}{RadGold},
    boxrule=.25pt,arc=0mm,left=5mm,right=4mm,top=3mm,bottom=3mm,#1}%
}{}

% Las once portadillas internas del propietario autónomo se convierten en
% divisiones científicas ordinarias dentro del capítulo 52.  Se evita así
% introducir partes autónomas, gigantismo tipográfico o páginas vacías.
\providecommand{\partopener}[3]{\section{#2}\noindent\emph{#3}\par\medskip}

\providecommand{\BigRadionEquation}{%
  \begin{equation*}
  \boxed{\displaystyle
  \frac{180^{\circ}}{\pih}
  =50^{\circ}+A^{\circ}-\frac{20}{81}\,\Delta_4^{\circ}+\epsR^{\circ}}
  \end{equation*}}
\makeatother


\chapter{El radión angular HMT--MD: empalme con acarreo de las secciones directa y torsional y cierre exacto de la unidad de fase, acción y memoria}
\label{chap:sucesor-102-radion}

\section{Composición APP--TRIT--TPK del estado de fase: generación de la unidad orientada con memoria}
% Fragmento público generado desde propietarios íntegros.
% Residencia: 52.1.
% Fuente: ../incoming/radion_monografia_20260827/manuscrito/sections/01_resumen_y_guia.tex:1-128
\paragraph{Síntesis matemática del objeto y de su cadena causal}

Esta monografía reconstruye el radián desde la cadena generativa
\[
\APP\longrightarrow\TRIT\longrightarrow\TPK
\longrightarrow X_{\mathrm{enr}}
\longrightarrow\text{estructura discreta del continuo}
\longrightarrow\MD,
\]
sin tomar como entrada ni el valor metrológico de \(180/\pi\), ni el pequeño
residuo angular que se desea explicar. El resultado es una unidad enriquecida,
el \emph{radión}, definida como la unidad orientada que recorre un radián de
fase y barre exactamente una acción \(\hret\) en la elipse positiva asociada
al mismo estado.

El cierre rector es
\BigRadionEquation
con
\[
S_E=2\pih\left(\frac5{36}+\frac{25}{9}\alphah\right),
\qquad
\Delta_4=
\frac{\pih-e_{\HMT}\log\pih}{270}
\left(1+\frac{\pih}{729}\right),
\]
\[
\rho_{\mathrm{tw}}=\frac{2\cdot90}{729}=\frac{20}{81},
\qquad
\epsone=\frac{20}{81}\Delta_4-(S_E-1).
\]
La división APP prueba \(1<S_E<2\), de modo que el cociente es exactamente
uno y el resto es \(D_E=S_E-1\). La igualdad
\[
1=S_E-\frac{20}{81}\Delta_4+\epsone
\]
es entonces una identidad interna. La carta angular posterior produce
\[
\epsR^\circ
=5.13830167585752820716851646226459474899085744\ldots
\times10^{-8}\,{}^\circ.
\]

El texto demuestra además que la elipse de acción
\[
Q=a_S\cos\theta,\qquad P=b_S\sin\theta,\qquad
a_Sb_S=2\hret
\]
satisface
\[
\operatorname{Area}(0,\theta)=\hret\theta.
\]
Por tanto, un radión barre \(\hret\) y la vuelta \(2\pi\) encierra
\(h_{\mathrm{ret}}=2\pi\hret\). La misma elipse tiene borde simpléctico
finito e interior de área infinita cuando se equipa con la métrica
Hilbert--Beltrami--Klein. Ésta es la forma exacta de la tesis física que guía
la obra: la superficie visible cierra la acción; la profundidad interior
conserva un desarrollo no agotable por la proyección plana.

La exposición mantiene tipadas cuatro realizaciones del mismo estado:

\begin{enumerate}[label=\textbf{\arabic*.}]
\item el círculo \(S^1\), que publica la fase;
\item la elipse positiva, que publica la acción y la anisotropía;
\item el interior de Klein, que mide la profundidad mediante razón doble;
\item la hélice nonádica, que conserva retorno, memoria y contracción.
\end{enumerate}

No son cuatro geometrías superpuestas sin criterio. Son cuatro lectores con
dominio, codominio y dato conservado diferentes, aplicados a un único estado
enriquecido. El sector \(90/120\), la dodecafase, el acarreo del radión, la cola
de Lambert, el núcleo de Barbero y las torres globales se presentan asimismo
como objetos emparentados pero no intercambiables.

\begin{thesisbox}
El radión es la unidad orientada que barre \(\hret\) en la elipse de acción.
El círculo publica su fase, Klein mide su profundidad, la hélice conserva su
memoria y el acarreo coinductivo empalma exactamente sus secciones directa y
torsional.
\end{thesisbox}

\paragraph{Orden expositivo y dependencias de lectura}

El orden no es ornamental. Las Partes I y II producen el objeto y su cierre
desde APP--TRIT--TPK. Sólo después se introducen las realizaciones en lenguaje
de Hilbert, LC, Maxwell, Minkowski, Hodge, Lorentz, modularidad y órbitas. Cada
capítulo distingue tres niveles:

\begin{description}[style=nextline,leftmargin=4.2cm,labelwidth=3.9cm]
\item[Resultado interno.] Igualdad o construcción demostrada en el dominio
HMT con sus primitivas declaradas.
\item[Formalización reunida.] Mapa nuevo que articula resultados ya existentes
sin usar el blanco como selector.
\item[Interfaz física.] Traducción dimensional posterior que reconoce una
estructura convencional sin hacerla origen de la selección.
\end{description}

Los recuadros de \emph{estatuto y falsador} indican qué observación concreta
invalidaría cada promoción fuerte. Esta disciplina no debilita la tesis; la
hace auditable. El objetivo de la narración extensa es que ninguna palabra
como \emph{círculo}, \emph{elipse}, \emph{torsión}, \emph{carga} o
\emph{memoria} oculte un cambio de tipo.

\paragraph{Resultados rectores y alcance probatorio}

\begin{enumerate}[label=\textbf{C\arabic*.}]
\item Cierre independiente del valor final del radión por dos publicaciones del mismo estado y
resta con acarreo APP.
\item Derivación del coeficiente \(20/81\) como dos hojas del sector activo
\(90\), normalizadas por la fibra \(3^6=729\).
\item Igualador tipado entre la costura \(\Re s=1/2\), la cara APP de cien
marcas y la fase dodecafásica \(\theta_2=50^\circ\).
\item Teorema de área del radión: \(\operatorname{Area}(\theta)=\hret\theta\).
\item Estudio focal exacto de la elipse y de sus invariantes euclídeos,
hiperbólicos, simplécticos y modulares.
\item Teorema borde--interior: acción total finita \(h\) e interior de Klein
de área hiperbólica infinita.
\item Realización exacta mediante covarianza mínima y oscilador LC, seguida de
la publicación electromagnética \(90/120\).
\item Atlas de no-identificaciones para todos los objetos \(90/120\), incluida
la separación entre \(\epsR\), cola espectral y núcleo Barbero \(12:-1\).
\item Caracterización estructural de \(\pi\) en la carta del radión, distinta
del generador coinductivo primario.
\item Síntesis ontológica: dimensión como rango mínimo de memorias que una
proyección de fase ya no puede conservar.
\end{enumerate}
% Fuente: ../incoming/radion_monografia_20260827/manuscrito/sections/01b_mapa_resultados.tex:1-133
\paragraph{Correspondencia entre resultados, tipos y residencias causales}

La correspondencia siguiente ordena las afirmaciones rectoras y sus demostraciones,
manteniendo visible la jerarquía que impide que
la riqueza física borre el cierre matemático que la sostiene.

% HMT_RESULT_BEGIN:RDN-01-GENEALOGIA:theorem
\begin{exactbox}[title={Genealogía APP--TRIT--TPK del radión}]
El radión nace de la composición efectiva
\(\APP\to\TRIT\to\TPK\to X_{\mathrm{enr}}\to\MD\); no se obtiene
redefiniendo retrospectivamente el radián convencional.
\end{exactbox}
% HMT_RESULT_END:RDN-01-GENEALOGIA:theorem

% HMT_RESULT_BEGIN:RDN-02-CIERRE:theorem
\begin{exactbox}[title={Cierre exacto e independencia prospectiva}]
El acarreo APP produce internamente
\[
1=S_E-\frac{20}{81}\Delta_4+\varepsilon_R^{(1)},
\qquad
\varepsilon_R^{(1)}=\frac{20}{81}\Delta_4-(S_E-1),
\]
y la carta angular transporta esta misma identidad a
\(180^\circ/\pi_{\!\HMT}\) sin introducir el residuo objetivo.
\end{exactbox}
% HMT_RESULT_END:RDN-02-CIERRE:theorem

% HMT_RESULT_BEGIN:RDN-03-BISAGRA-50:theorem
\begin{exactbox}[title={Coordenada crítica de \(50^\circ\)}]
La coordenada crítica \(\Re s=1/2\), publicada en la cara APP de cien
marcas y en la rueda dodecafásica, selecciona la fase
\(\theta_2=50^\circ\). Es una igualdad de lectores tipados del mismo estado.
\end{exactbox}
% HMT_RESULT_END:RDN-03-BISAGRA-50:theorem

% HMT_RESULT_BEGIN:RDN-04-COEFICIENTE-2081:theorem
\begin{exactbox}[title={Densidad bidireccional nonádica \(20/81\)}]
El coeficiente torsional no se ajusta en la red \(1/81\): se construye como
\(\rho_{\mathrm{tw}}=2N_6/3^6=2\cdot90/729=20/81\), conservando hoja,
orientación y normalización de la fibra nonádica.
\end{exactbox}
% HMT_RESULT_END:RDN-04-COEFICIENTE-2081:theorem

% HMT_RESULT_BEGIN:RDN-05-AREA-ACCION:theorem
\begin{exactbox}[title={Ley areal del radión}]
Para la elipse positiva
\(Q=a_S\cos\theta\), \(P=b_S\sin\theta\),
\(a_Sb_S=2\hbar_{\mathrm{ret}}\), el área orientada barrida es
\(\mathcal A(\theta)=\hbar_{\mathrm{ret}}\theta\). Un radión barre exactamente
\(\hbar_{\mathrm{ret}}\), y la vuelta \(2\pi\) encierra
\(h_{\mathrm{ret}}=2\pi\hbar_{\mathrm{ret}}\).
\end{exactbox}
% HMT_RESULT_END:RDN-05-AREA-ACCION:theorem

% HMT_RESULT_BEGIN:RDN-06-INVARIANTES-FOCALES:theorem
\begin{exactbox}[title={Invariantes focales del operador positivo}]
El par publicado \(A\pm C^\ast\) fija los semiejes, la excentricidad, la
separación focal y el parámetro de compresión simpléctica. En particular,
\((d_1/d_0)^2=C^\ast/2\), identidad adimensional que separa forma de escala.
\end{exactbox}
% HMT_RESULT_END:RDN-06-INVARIANTES-FOCALES:theorem

% HMT_RESULT_BEGIN:RDN-07-BORDE-INTERIOR:theorem
\begin{exactbox}[title={Finitud simpléctica del borde y divergencia del área hiperbólica interior}]
El borde simpléctico de la elipse encierra la acción finita
\(h_{\mathrm{ret}}\), mientras el área Hilbert--Beltrami--Klein del interior
diverge al alcanzar la frontera. Fase visible y profundidad no son la misma
medida.
\end{exactbox}
% HMT_RESULT_END:RDN-07-BORDE-INTERIOR:theorem

% HMT_RESULT_BEGIN:RDN-08-HELICE-MEMORIA:theorem
\begin{exactbox}[title={Elevación helicoidal nonádica con retorno de fase y avance de memoria}]
La monodromía nonádica devuelve la fase sin reiniciar el estado. La hélice
levanta el círculo conservando acarreo y memoria; el cierre \(2\pi/4\pi\)
distingue retorno proyectivo, inversión de orientación y retorno completo.
\end{exactbox}
% HMT_RESULT_END:RDN-08-HELICE-MEMORIA:theorem

% HMT_RESULT_BEGIN:RDN-09-OSCILADOR-EM:development
\begin{narrativebox}[title={Realización oscilatoria y constitutiva electromagnética}]
La misma anisotropía se realiza mediante una covarianza mínima y un oscilador
LC: una coordenada almacena el aspecto eléctrico-elástico y su conjugada el
aspecto magnético-rotacional. La frecuencia y la impedancia quedan tipadas por
separado.
\end{narrativebox}
% HMT_RESULT_END:RDN-09-OSCILADOR-EM:development

% HMT_RESULT_BEGIN:RDN-10-TIPOS-90120:theorem
\begin{exactbox}[title={Separación tipada de los operadores \(90/120\)}]
Extractor dodecafásico, acarreo del radión, cola de Lambert, funcional
hexada--octada, núcleo de Barbero y semigrupo global de torres son objetos
distintos. El peso \(12:-1\) procede de los doce sectores de la cadena
fundamental dodecafásica y de su única costura orientada. El coeficiente de
\((1-q)^{-1}\), igual a \(1/24\), y el contratérmino entero mínimo verifican
posteriormente el par producido.
\end{exactbox}
% HMT_RESULT_END:RDN-10-TIPOS-90120:theorem

% HMT_RESULT_BEGIN:RDN-11-GEOMETRIAS-CAMPO:development
\begin{narrativebox}[title={Realizaciones compleja, lorentziana y dual de Hodge}]
Los regímenes elíptico, parabólico e hiperbólico del TRIT se publican como
rotación, umbral y transformación hiperbólica. La estrella de Hodge, las constitutivas y la fuerza
de Lorentz actúan después sobre el estado ya orientado; no seleccionan sus
constantes.
\end{narrativebox}
% HMT_RESULT_END:RDN-11-GEOMETRIAS-CAMPO:development

% HMT_RESULT_BEGIN:RDN-12-MODULAR-ORBITAL:development
\begin{narrativebox}[title={Caracteres modular, orbital y aritmético de frontera}]
La involución modular, las órbitas elípticas, el avance-retardo y la aritmética
3--6--9 reconocen distintas memorias del mismo estado. La función modular y
las vacancias se usan como lectores posteriores, nunca como sustitutos del
generador APP--TRIT--TPK.
\end{narrativebox}
% HMT_RESULT_END:RDN-12-MODULAR-ORBITAL:development

% HMT_RESULT_BEGIN:RDN-13-SINTESIS-ONTOLOGICA:conclusion
\begin{thesisbox}[title={Síntesis operatoria de fase, acción, profundidad y memoria}]
El círculo es la proyección de fase; la elipse es la sección de acción; Klein
mide la profundidad; la hélice conserva la historia. El radión es la unidad
orientada que mantiene acopladas esas publicaciones sin borrar el residuo que
las distingue.
\end{thesisbox}
% HMT_RESULT_END:RDN-13-SINTESIS-ONTOLOGICA:conclusion

% HMT_RESULT_BEGIN:RDN-14-CONTROL-ALCANCE:control
\begin{statusbox}[title={Delimitación de dominios y continuidad causal}]
Las leyes completas de masas, el refinamiento electrónico, CKM/CP y el
corredor excepcional conservan sus demostraciones rectoras propias. Esta monografía los
cita sólo donde comparten antecedente o carta; no traslada al radión sus
selectores específicos ni confunde sus residencias probatorias.
\end{statusbox}
% HMT_RESULT_END:RDN-14-CONTROL-ALCANCE:control


\subsection{Soporte bivariado de APP, orientación ternaria del TRIT y actualización del estado enriquecido mediante el TPK}
% Fragmento público generado desde propietarios íntegros.
% Residencia: 52.1.1.
% Fuente: ../incoming/radion_monografia_20260827/manuscrito/sections/02_genealogia.tex:1-180
\noindent La genealogía operatoria antecede a la medida angular y determina las variables conservadas por cada lector.\par\medskip

\subsubsection{Soporte APP--TRIT--TPK y construcción del estado enriquecido}

\paragraph{Antecedencia del estado enriquecido respecto de la publicación circular}

La definición escolar de radián comienza con un círculo ya dado, un centro ya
dado y una longitud de arco ya mensurable. Si el radio y el arco tienen la
misma longitud, el ángulo subtendido mide un radián. La definición es correcta
dentro de la geometría euclídea, pero no responde a una pregunta anterior:
¿qué estructura hace posible que una vuelta posea fase, orientación, acción y
memoria, y por qué la unidad natural se relaciona con \(2\pi\)?

HMT invierte el orden causal. Primero construye el estado finito, después su
prolongación coinductiva, luego la vuelta y sólo al final su publicación
angular. El radión no corrige el radián convencional; lo levanta a su dominio
genealógico.

\begin{typebox}
\[
\APP\to\TRIT\to\TPK\to X_{\mathrm{enr}}
\to\mathcal C^{\mathrm{disc}}_{\mathrm{cont}}
\to\bigl(\pih,e_{\HMT},\varphi_{\HMT},\alphah,A,C^*,\Delta_4,\hret\bigr)
\to\mathfrak r_\sigma.
\]
La matemática convencional sólo aparece después como lenguaje de realización:
\(S^1\), elipse simpléctica, métrica de Klein, espacio de Hilbert, circuito
LC, pantalla de Minkowski y ecuaciones de campo.
\end{typebox}

\paragraph{Soporte bivariado de APP: hojas, residuos, cocientes y acarreo}

La Plataforma de Aritmética Posicional (APP) conserva simultáneamente dos
realizaciones de los símbolos \(1,\dots,9\): una hoja aditiva y una hoja
multiplicativa. No se reduce un número a su valor publicado. Junto al valor
permanecen hoja, orientación, cociente, resto, acarreo, frontera, ruta y
supervivencia.

El censo básico tiene tres estratos que jamás deben identificarse por mera
coincidencia cardinal:
\[
104,976\longrightarrow468\;U_6\longrightarrow243\;w_6
\subset\F_3^6,\qquad |\F_3^6|=729.
\]
Las \(468\) emisiones decimales, las \(243\) palabras visibles y el ambiente
de \(729\) palabras son objetos distintos. El ambiente será decisivo para la
normalización \(20/81\).

En base cúbica \(B=1000\), la resta orientada de dos palabras de tríadas se
realiza de derecha a izquierda:
\begin{equation}
u_i=t_i-v_i+c_{i+1},\qquad
r_i=u_i\bmod1000,\qquad
c_i=\frac{u_i-r_i}{1000}.
\label{eq:subacarreo}
\end{equation}
Fijados los operandos y el acarreo remoto, cociente y resto son únicos. Éste es
el operador que, más adelante, publicará las cifras de \(\epsone\). No hay
ningún decimal elegido después de ver el radián.

La completación
\[
10^3=729+243+27+1
\]
no es una curiosidad decimal. Expresa la conservación de la unidad al pasar
del ambiente ternario a una cámara de mil posiciones. La diferencia
\(1000-729=271\) es la corona completante; la identidad correcta es
\(999=729+270\), no \(729+271\).

\paragraph{Orientación ternaria y clasificación de los 3 regímenes operatorios del TRIT}

El TRIT asigna a cada estado local uno de tres regímenes orientados. En la
convención temporal activa,
\[
\begin{aligned}
+1&\longleftrightarrow\text{retorno, memoria, pasado},\\
0&\longleftrightarrow\text{umbral, presente},\\
-1&\longleftrightarrow\text{apertura bidireccional, futuro}.
\end{aligned}
\]
La realización cuadrática distingue:
\begin{equation}
J^2=-I,\qquad N^2=0,\qquad R^2=I.
\label{eq:tres-regimenes}
\end{equation}
Así aparecen, respectivamente, la rotación elíptica, el umbral parabólico y
la separación hiperbólica. No se mezclan tres métricas. Se conservan tres
acciones tipadas sobre un antecedente común.

La raíz interna de \(-1\) no se importa suponiendo de antemano el plano
complejo. En el bloque APP, dos proyecciones \(P,Q\) con razón
\(r=3/4=90/120\) producen
\[
J_{\mathrm{comm}}=\frac{[P,Q]}{\sqrt{r(1-r)}},\qquad J_{\mathrm{comm}}^2=-I,
\]
mientras
\[
R_{\mathrm{comm}}=2P-I,\qquad R_{\mathrm{comm}}^2=I,\qquad
R_{\mathrm{comm}}J_{\mathrm{comm}}=-J_{\mathrm{comm}}R_{\mathrm{comm}}.
\]
La pareja genera la forma mínima de \(\mathrm{Cl}_{1,1}\simeq M_2(\R)\):
rotación y apertura emergen juntas, pero siguen siendo operaciones diferentes.

\paragraph{Composición dinámica del TPK: selección, transporte, memoria y retorno}

El Topological Phase Núcleo no es un nombre para una caja de objetos. Es la
composición que selecciona fase, transporta cardinalidad, levanta hojas,
registra memoria y ejecuta el retorno. La dependencia fundamental es
\[
U_t\longrightarrow\Gamma_9\longrightarrow K_{\mathrm{ph}}.
\]
La holonomía nonádica \(\Gamma_9\) actúa sobre el estado enriquecido; la
publicación \(K_{\mathrm{ph}}\) es posterior y conserva memoria reducida. La
identidad \(K_{\mathrm{ph}}^9|_r=E_+E_\times\) no convierte a \(K_{\mathrm{ph}}\) en
generador.

El certificado nonádico conserva la cadena completa
\[
w_6\to w_{12}\to w_{18}\to w_{24}\to w_{30}
\to R_{36}\to G_9,
\]
con \(396\) niveles, \(2376\) trits, \(43\) vueltas, \(387\) pares
\((K,K+9)\) por canal y cinco regiones de \(\pi\). La fase cumple
\[
g(k+9)=g(k),
\]
pero el estado completo no retorna:
\[
B_{k+9}\ne B_k.
\]
Ésta es la semilla matemática de la hélice: la proyección angular cierra; la
memoria avanza.

\paragraph{Extensión temporal no escindida \(C_{12}\to C_{108}\to C_9\)}

La memoria de doce pasos se organiza mediante
\begin{equation}
0\longrightarrow C_{12}\longrightarrow C_{108}
\longrightarrow C_9\longrightarrow0.
\label{eq:extension-108}
\end{equation}
El cociclo
\[
c(r,r')=\left[\left\lfloor\frac{r+r'}9\right\rfloor\right]_{12}
\]
mide el acarreo entre retorno de fase y avance de memoria. Nueve pasos cierran
la fase observable \(C_9\); no reinician el registro dodecafásico.

Conviene separar cuatro objetos:
\begin{center}
\begin{tabularx}{.95\textwidth}{>{\bfseries}l X}
\toprule
\(C_{12}\) & grupo cíclico de doce posiciones;\\
\(\Rstate\) & registro enriquecido \((E_m,\tau_m,\sigma_m,\kappa_m,\Lambda_m)_{m=1}^{12}\);\\
\(\Theta_{108}\) & publicación cruda en el calendario de 108;\\
\(\Gamma_{108}\) & historia u holonomía levantada.\\
\bottomrule
\end{tabularx}
\end{center}
Por ello, la expresión vaga \enquote{post-\(C_{12}\)} se sustituirá por
\enquote{posterior al registro dodecafásico \(\Rstate\), en su lectura
sexádica} cuando ése sea el mapa efectivo.

\paragraph{Emergencia correlativa de las 5 construcciones del continuo desde el estado único}

La estructura discreta del continuo no se obtiene ensamblando cinco teorías
independientes. Sus cinco aspectos internos ---espectral, solenoidal, calibre,
cohomológico y determinantal--- emergen correlativamente del mismo estado
APP--TRIT--TPK. Esta obra no vuelve a demostrar el tratado integral, pero
preserva el hecho causal que necesita el radión: la sección coinductiva que
genera \(\pi,e,\varphi\) y \(\alpha\) no vive en una recta real desnuda; vive
en un objeto con fibras, relaciones, orientación, acarreo, memoria, terminal y
lector.

\begin{statusbox}
\textbf{Falsador de genealogía.} Si una fórmula posterior recibe
\(180/\pi\), \(\epsR\), una masa o una tabla metrológica para elegir una fase,
un coeficiente o una rama, deja de ser una generación HMT-first. El contraste
externo puede reconocer una salida; no puede seleccionarla.
\end{statusbox}


\subsection{Independencia prospectiva del cierre respecto del valor final}
% Fragmento público generado desde propietarios íntegros.
% Residencia: 52.1.2.
% Fuente: ../incoming/radion_monografia_20260827/manuscrito/sections/02_genealogia.tex:181-272

\subsubsection{Genealogía de las constantes y construcción del par angular}

\paragraph{Sección coinductiva y unicidad del valor generado de \(\pi\)}

Sea \(x_\infty\) la sección global del sistema inverso. Los lectores
coinductivos producen
\[
\operatorname{ev}_\pi(x_\infty)=\pih,\qquad
\operatorname{ev}_e(x_\infty)=e_{\HMT},\qquad
\operatorname{ev}_\varphi(x_\infty)=\varphi_{\HMT}.
\]
El subíndice HMT no designa otro número. Conserva la genealogía de la misma
constante. En particular,
\[
\pih=3.14159265358979323846264338327950288419716939937510\ldots.
\]
La vuelta positiva es
\begin{equation}
T_+(x_\infty)=2\pih.
\label{eq:vuelta-plus}
\end{equation}
Su publicación angular de una unidad es
\begin{equation}
R_{\HMT}^{\circ}=\frac{360^\circ}{T_+}=\frac{180^\circ}{\pih}.
\label{eq:radian-publicado}
\end{equation}
La ecuación no define todavía el radión enriquecido; define la carta visible
del radián una vez generada la vuelta.

\paragraph{Bidireccionalidad de rutas y prolongaciones conjugadas}

La bidireccionalidad no se enuncia como eslogan. La involución concreta
\begin{equation}
J_\alpha(T)=\frac{\alphah^2}{T},\qquad
T_-=J_\alpha(T_+),\qquad T_+T_-=\alphah^2
\label{eq:bidireccional-turn}
\end{equation}
conserva un producto y revierte la hoja. Expansión y contracción son dos
prolongaciones conjugadas del mismo estado. La raíz finita de un jet
cuadrático puede aproximar \(T_+\), pero no sustituye a la sección
coinductiva: la diferencia certificada es
\[
T_{P_9}-T_+=-5.1110566290335\ldots\times10^{-25}.
\]
El control finito es un testigo; el fundamento es el límite.

\paragraph{Dependencia causal entre \(\alpha\), acción y publicación angular}

El cierre dodecafásico genera \(\alphah\) desde el estado firmado y sus
cartas reversibles. Después, el lector determinantal local fija la década de
acción mediante
\[
\Sigma_H=\varphi_{\HMT}\alphah^{16},
\qquad \operatorname{ord}_{10}\Sigma_H=-34.
\]
La sección de retorno produce \(\hret\). Ni \(\alpha\) ni \(\hbar\) se
\enquote{miden en grados}. Se publican posteriormente mediante el par angular
\begin{equation}
P(\alpha,\eta_{\mathrm{ret}})=
\left([10^3\alpha]_{360},[2(\eta_{\mathrm{ret}}+\alpha)]_{360}\right)
=(A,C^*).
\label{eq:parangular}
\end{equation}
En la rama canónica,
\[
A=7.297352569283800997285105472380663\ldots{}^\circ,
\]
\[
C^*=2.1237383389686457242619513803933607937\ldots{}^\circ.
\]
Sólo entonces la carta analítica
\[
x=\frac{\pi A}{180},\qquad y=\frac{\pi C^*}{180}
\]
abre los canales \(q_\pm=e^{-x\mp y}\).

El orden causal completo es
\[
\alpha\to(\eta\parallel-34)\to\hbar
\to P\to\kappa_{360}\to q_\pm
\to\text{realizaciones MD}.
\]
La relación \(C^*/2-\alpha\) puede servir de verificación inversa. No se usa
para generar \(\hret\) retroactivamente.

\begin{narrativebox}
La carta angular no convierte una constante adimensional o una acción en un
ángulo por decreto. Publica dos coordenadas del estado en una rueda finita.
El grado es el idioma de esa publicación; no es la sustancia ontológica de
\(\alpha\) ni de \(\hbar\).
\end{narrativebox}
% Fuente: ../incoming/radion_monografia_20260827/manuscrito/sections/03_cierre_exacto.tex:254-285
\paragraph{Independencia prospectiva respecto del valor final}

El verificador del operador de profundidad prohíbe expresamente tres entradas:
\[
\frac{180}{\pi},\qquad \epsR,\qquad
\text{el decimal objetivo}.
\]
Primero construye \(S_E\) y \(S_T\), después ejecuta la resta APP, y sólo al
final compara el cierre. Los controles dan
\[
\texttt{target\_epsilon\_used=False},\qquad
\texttt{target\_180\_over\_pi\_used=False}.
\]
Asimismo:
\begin{itemize}
\item una sola hoja da un defecto de orden \(10^{-5}\);
\item sustituir \(\theta_2\) por \(\theta_1\) o \(\theta_3\) desplaza el
resultado en aproximadamente \(\pm\pi/6\);
\item eliminar el canal \(e\) deja un defecto de orden \(10^{-3}\);
\item la variante histórica \(\Delta_4(A,C^*)\) produce otro número;
\item la cola espectral posterior a \(\Rstate\) y el operador armónico de
cuatro términos tampoco coinciden con \(\epsone\).
\end{itemize}

\begin{statusbox}
\textbf{Estatuto.} El cierre es una \status{formalización nueva} demostrada
con estructura de partida explícita y certificado computacional independiente del valor final.
La igualdad es exacta en el límite coinductivo. El perfil corto de 36 cifras
de \(\alpha\) genera una variante que diverge sólo después de muchas cifras;
se conserva como testigo histórico, no como valor canónico.
\end{statusbox}



\section{Secciones directa y torsional: acarreo APP y cierre exacto del radión}
\subsection{Construcción independiente de las secciones \(S_E\) y \(S_T\) antes de definir el cociclo \(\varepsilon_R=S_T-S_E\)}
% Fragmento público generado desde propietarios íntegros.
% Residencia: 52.2.1.
% Fuente: ../incoming/radion_monografia_20260827/manuscrito/sections/03_cierre_exacto.tex:1-50
\noindent El cierre exacto se obtiene mediante dos secciones construidas antes del cociclo que las relaciona.\par\medskip

\subsubsection{Construcción independiente de las secciones directa y torsional}

\paragraph{Sección directa del sistema temporal dodecafásico}

La rueda dodecafásica tiene centros
\begin{equation}
\theta_m=30m-10^\circ,\qquad m=1,\ldots,12.
\label{eq:rueda-dodecafasica}
\end{equation}
La segunda fase es \(\theta_2=50^\circ\). Junto a la publicación
\(A=10^3\alphah\), forma la sección directa
\begin{align}
S_E
&=\frac{T_+}{360}\,(50+10^3\alphah)\\
&=2\pih\left(\frac5{36}+\frac{25}{9}\alphah\right).
\label{eq:seccion-electrica}
\end{align}
El subíndice \(E\) significa aquí \emph{directa o común}; su realización
eléctrica llegará después del lector constitutivo. No se postula
literalmente \(A=E\).

Los cilindros coinductivos prueban
\begin{equation}
1<S_E<2.
\label{eq:SE-intervalo}
\end{equation}
La división APP fuerza entonces
\begin{equation}
q_E=\lfloor S_E\rfloor=1,
\qquad
D_E=\operatorname{rem}_1(S_E)=S_E-1.
\label{eq:DE}
\end{equation}
El entero uno no es el número de vueltas de la hélice. Es el cociente único
de la sección directa dentro de la celda de unidad.

\paragraph{Sección torsional determinada por el operador \(\Delta_4\)}

La segunda sección utiliza una salida coinductiva anterior:
\begin{equation}
\Delta_4=
\frac{\pih-e_{\HMT}\log\pih}{270}
\left(1+\frac{\pih}{729}\right).
\label{eq:delta4}
\end{equation}
La palabra \emph{torsión} designa aquí el cuanto global de la carta HMT. No
lo identifica automáticamente con la torsión de Cartan
\(T^I=D_\omega e^I\). El puente hacia Cartan exige otra flecha tipada.

\subsection{Densidad orientada del sector activo y determinación única de \(20/81\)}
% Fragmento público generado desde propietarios íntegros.
% Residencia: 52.2.2.
% Fuente: ../incoming/radion_monografia_20260827/manuscrito/sections/03_cierre_exacto.tex:51-103

\paragraph{Densidad orientada \(20/81\) del sector activo \(90/120\)}

El selector trítico de fase separa
\[
H_+=\{1,4,7\},\qquad H_-=\{2,5,8\},\qquad H_0=\{3,6,9\}.
\]
Hay seis fases activas y quince parejas no ordenadas:
\[
H_{\mathrm{act}}=H_+\cup H_-,\qquad
\left|\binom{H_{\mathrm{act}}}2\right|=\binom62=15.
\]
El modo activo y su extensión a ocho posiciones son
\begin{equation}
N_6=6\binom62=90,\qquad
N_8=8\binom62=120.
\label{eq:90-120-incidencia}
\end{equation}
Dos hojas conjugadas del sector activo, normalizadas por la fibra ambiente,
dan
\begin{equation}
\rho_{\mathrm{tw}}
=\frac{2N_6}{|\F_3^6|}
=\frac{2\cdot90}{729}
=\frac{180}{729}
=\boxed{\frac{20}{81}}.
\label{eq:rho-twoway}
\end{equation}

\begin{proposition}[Unicidad relativa de la densidad]
Fijados el sector activo \(N_6=90\), las dos hojas conjugadas y la
normalización por \(|\F_3^6|=729\), la densidad orientada es necesariamente
\(20/81\).
\end{proposition}
\begin{proof}
La cardinalidad transportada es \(2N_6=180\). La normalización por el ambiente
da \(180/729\); dividir numerador y denominador por \(9\) produce \(20/81\).
No interviene ninguna comparación con \(180/\pi\) ni con \(\epsR\).
\end{proof}

La lectura histórica como proyección sobre una red \(1/81\) es un control de
rigidez útil, pero ya no es el fundamento: el coeficiente se obtiene antes
del cierre. La ablación de una hoja produce \(10/81\), y la salida cambia en
orden \(10^{-5}\); por tanto, el factor dos no es decorativo.

La sección torsional es
\begin{equation}
\Phi_T=\rho_{\mathrm{tw}}\Delta_4,
\qquad
S_T=1+\Phi_T.
\label{eq:seccion-torsional}
\end{equation}


\subsection{El acarreo \(\varepsilon_R\) como cociclo de transición entre secciones}
% Fragmento público generado desde propietarios íntegros.
% Residencia: 52.2.3.
% Fuente: ../incoming/radion_monografia_20260827/manuscrito/sections/03_cierre_exacto.tex:104-253
\subsubsection{Cociclo de transición entre las secciones directa y torsional}

\paragraph{Definición operatoria del acarreo \(\varepsilon_R\)}

\begin{definition}[Acarreo real del radión]
La salida unitaria del radión es el cociclo de transición
\begin{equation}
\epsone:=S_T-S_E
=\Phi_T-D_E
=\frac{20}{81}\Delta_4-operatorname{rem}_1(S_E).
\label{eq:epsilon-def}
\end{equation}
\end{definition}

La ecuación contiene una resta, pero no es un ajuste al blanco. Sus dos
operandos \(S_T\) y \(S_E\) se generan separadamente. La independencia que
se demuestra es independencia respecto de \(180/\pi\) y de un
\(\varepsilon\) objetivo; no se afirma independencia algebraica entre una
diferencia y sus sumandos.

\begin{theorem}[Cierre exacto unitario]
Con \(S_E\), \(\Phi_T\) y \(\epsone\) definidos por
\eqref{eq:seccion-electrica}, \eqref{eq:delta4}, \eqref{eq:rho-twoway} y
\eqref{eq:epsilon-def}, se cumple exactamente
\begin{equation}
\boxed{1=S_E-\Phi_T+\epsone.}
\label{eq:cierre-unitario}
\end{equation}
\end{theorem}
\begin{proof}
Por \eqref{eq:SE-intervalo}, \(D_E=S_E-1\). Por definición,
\(\epsone=\Phi_T-D_E\). Entonces
\[
S_E-\Phi_T+\epsone
=S_E-\Phi_T+\Phi_T-(S_E-1)=1.
\]
La fuerza de la prueba no está en esta cancelación elemental, sino en que la
genealogía previa fija los dos operandos sin usar el cierre como entrada.
\end{proof}

\paragraph{Publicación del cierre en la carta angular}

La carta de la vuelta multiplica una fracción unitaria por
\(360/T_+=180/\pih\). En particular,
\[
\epsR^\circ=\frac{360^\circ}{T_+}\epsone.
\]
Multiplicar \eqref{eq:cierre-unitario} por \(360^\circ/T_+\) y usar
\[
\frac{360}{T_+}S_E
=360\left(\frac5{36}+\frac{25}{9}\alphah\right)
=50+1000\alphah=50+A
\]
produce el cierre angular.

\begin{theorem}[Ecuación exacta del radión]
\BigRadionEquation
\end{theorem}

La ecuación dice algo más preciso que \enquote{el radián es
\(57.2957\ldots{}^\circ\)}. Descompone esa publicación en cuatro operaciones
con demostración de referencia:
\begin{center}
\begin{tabularx}{.96\textwidth}{>{\bfseries}l X}
\toprule
\(50^\circ\) & fase segunda de la rueda; bisagra crítica en una carta declarada;\\
\(A^\circ\) & coordenada común de la publicación parangular;\\
\(-\frac{20}{81}\Delta_4^\circ\) & transporte orientado de la torsión global;\\
\(+\epsR^\circ\) & acarreo real que empalma los dos lifts.\\
\bottomrule
\end{tabularx}
\end{center}

\paragraph{Convergencia por profundidad del refinamiento}

A profundidad \(N\), el generador entrega cilindros racionales
\(p_N\to\pih\) y \(e_N\to e_{\HMT}\). Se definen
\[
S_{E,N}=2p_N\left(\frac5{36}+\frac{25}{9}\alphah\right),
\]
\[
\Delta_{4,N}=\frac{p_N-e_N\log p_N}{270}
\left(1+\frac{p_N}{729}\right),\qquad
S_{T,N}=1+\frac{20}{81}\Delta_{4,N},
\]
y
\[
\mathcal E_N=S_{T,N}-S_{E,N}.
\]
La continuidad del funcional en \(3<p<4\), \(2<e<3\), junto a las
derivadas monótonas, transporta cada par de cilindros a un intervalo
certificado. A profundidad \(45\) bloques, su anchura es menor que
\(4.81\times10^{-130}\).

Cada retorno de nueve niveles contiene seis trits por nivel. Por eso la tasa
de refinamiento es
\begin{equation}
729^{-9}=3^{-54}.
\label{eq:contraccion-54}
\end{equation}
Esta cifra controla la anchura de los intervalos; no es el valor de
\(\epsR\).

La resta de tríadas estabiliza
\begin{center}
\ttfamily
000|000|000|896|802|822|044|562|981|999|050|695|020|651|286|413
\end{center}
y el prefijo de carries
\begin{center}
\ttfamily
0,0,0,0,0,-1,0,-1,-1,-1,0,-1,0.
\end{center}
Éste es el mecanismo concreto de memoria decimal: no una cola elegida, sino
la recurrencia \eqref{eq:subacarreo} aplicada a operandos previamente
construidos.

\paragraph{Evaluación de los valores canónicos}

La salida coinductiva es
\begin{align}
\epsone
&=8.9680282204456298199905069502065128641372736205009\ldots
\times10^{-10},\\
\epsR^\circ
&=5.1383016758575282071685164622645947489908574400054\ldots
\times10^{-8}\,{}^\circ.
\label{eq:epsilon-valores}
\end{align}
La publicación completa resulta
\[
\frac{180^\circ}{\pih}
=57.2957795130823208767981548141051703324054724665643\ldots{}^\circ.
\]

\begin{exactbox}
\begin{align*}
50^\circ+A^\circ
&=57.29735256928380099728510547238066299956\ldots{}^\circ,\\
\frac{20}{81}\Delta_4^\circ
&=0.00157310758449687906223272996065728980465\ldots{}^\circ,\\
50^\circ+A^\circ-\frac{20}{81}\Delta_4^\circ
&=57.29577946169930411822287274242000570976\ldots{}^\circ,\\
\epsR^\circ
&=0.00000005138301675857528207168516462264594749\ldots{}^\circ,\\
\text{suma final}
&=57.29577951308232087679815481410517033241\ldots{}^\circ.
\end{align*}
\end{exactbox}



\section{Coordenada crítica y sistema temporal dodecafásico orientado: publicación afín de la fase de \(50^\circ\)}
\subsection{Aplicación afín \(j_{100}\) de la carta APP y unicidad de la fase \(m=2\) en \(j_{100}(1/2)=\theta_m=50^\circ\)}
% Fragmento público generado desde propietarios íntegros.
% Residencia: 52.3.1.
% Fuente: ../incoming/radion_monografia_20260827/manuscrito/sections/04_cincuenta_critica.tex:1-105
\noindent La coordenada crítica enlaza la carta de Cayley, el registro APP y la fase dodecafásica mediante aplicaciones tipadas.\par\medskip

\subsubsection{Dos realizaciones exactas de la coordenada angular \(50^\circ\)}

\paragraph{Segunda fase del sistema temporal dodecafásico}

Por \eqref{eq:rueda-dodecafasica}, la sucesión de centros comienza
\[
20^\circ,\ 50^\circ,\ 80^\circ,\ 110^\circ,\ldots.
\]
Resolver
\[
30m-10=50
\]
da únicamente \(m=2\). Por tanto, \(50^\circ\) no es una redondeo de \(A\),
una media vuelta ni el centro de un círculo: es una fase tipada del calendario
dodecafásico. La fracción de vuelta es
\[
\frac{50}{360}=\frac5{36},
\]
exactamente el primer sumando de \(S_E/T_+\).

\paragraph{Igualación de la coordenada crítica de Cayley}

La involución funcional
\[
\rho\longmapsto1-\bar\rho
\]
fija
\[
\operatorname{Fix}(\rho\mapsto1-\bar\rho)
=\{\rho:\Re\rho=1/2\}.
\]
El número \(1/2\) es la coordenada real autodual del intervalo funcional
\(0\leq\Re s\leq1\). No es \enquote{la mitad de una recta} en sentido
longitudinal, y no contiene la altura espectral \(\gamma\). Para
\[
\rho_\gamma=\frac12+i\gamma,
\]
la transformación de Cayley se puede escribir
\begin{equation}
\tau_\gamma
=\frac{\rho_\gamma-1}{\rho_\gamma}
=\frac{\gamma+i/2}{\gamma-i/2}
\in S^1,
\qquad
\gamma=\frac12\cot\frac{\theta_\gamma}{2}.
\label{eq:cayley-critica}
\end{equation}
Así, \(1/2\) fija la costura, \(\gamma\) determina el carácter angular y
\(\Gamma_9\) conserva la profundidad de sus iteraciones.

\paragraph{Inyección afín \(j_{100}\) del registro APP en la carta angular}

La completación cúbica dispone una cámara de diez marcas por coordenada. Una
cara contiene \(10^2=100\) posiciones. Declaramos la carta afín
\begin{equation}
j_{100}:[0,1]\longrightarrow[0,100^\circ],
\qquad j_{100}(\sigma)=100\sigma^\circ.
\label{eq:j100}
\end{equation}
Entonces
\[
j_{100}\left(\frac12\right)=50^\circ.
\]
Al exigir simultáneamente \(j_{100}(1/2)=\theta_m\), resulta
\[
100\cdot\frac12=30m-10
\quad\Longleftrightarrow\quad m=2.
\]

\begin{theorem}[Igualador crítico--dodecafásico]
En la carta APP de cien marcas, la coordenada fija de la involución crítica
se publica exactamente en la segunda fase de la rueda:
\begin{equation}
\boxed{j_{100}(1/2)=\theta_2=50^\circ.}
\label{eq:igualador-50}
\end{equation}
La fase \(m=2\) es única.
\end{theorem}

El teorema es exacto una vez declarada \eqref{eq:j100}. Su estatuto es
\status{formalización nueva}: el medio crítico, la completación \(10^3\), la
cara de cien marcas y \(\theta_2\) ya estaban construidos; la elección de esa
cara como carta crítica afín reúne por primera vez sus dominios.

\paragraph{Compatibilidad del valor crítico con la carta de Cayley}

La transformación de Cayley no envía \(1/2\) a \(50^\circ\). Si
\(\gamma=0\), \eqref{eq:cayley-critica} da
\[
\tau_0=-1=e^{i\pi},
\]
es decir, \(180^\circ\). Los \(50^\circ\) pertenecen a la carta APP de cien
marcas y al calendario dodecafásico; el ángulo Cayley depende de \(\gamma\).
Esta distinción impide colapsar todos los ceros de la función zeta en una sola
fase.

\begin{falsifier}
La identificación \(1/2\leftrightarrow50^\circ\) falla si se elimina la
carta \eqref{eq:j100}. La igualdad escalar \(100(1/2)=50\) no basta para
identificar por sí sola dos objetos. El teorema depende de la cara APP
distinguida y de la rueda dodecafásica.
\end{falsifier}


\subsection{Compatibilidad de la fase crítica con el retorno nonádico y el avance de memoria}
% Fragmento público generado desde propietarios íntegros.
% Residencia: 52.3.2.
% Fuente: ../incoming/radion_monografia_20260827/manuscrito/sections/04_cincuenta_critica.tex:106-171
\subsubsection{Separación tipada de los 2 objetos de valor \(1/2\)}

APP contiene además un modo singular con
\[
s=\frac12,\qquad s^2=\frac14,\qquad
1-s^2=\frac34=\frac{90}{120},\qquad
s\sqrt{1-s^2}=\frac{\sqrt3}{4}.
\]
Este \(s\) nace de los ángulos principales de las hojas suma y producto. La
costura crítica \(\Re\rho=1/2\) nace de una involución funcional. Que ambas
coordenadas valgan \(1/2\) es una coincidencia escalar con posible contenido,
pero no autoriza a identificarlas sin un cuadrado tipado.

La forma correcta de conservar el hallazgo es un diagrama de igualadores:
\[
\begin{tikzpicture}[baseline=(current bounding box.center),node distance=14mm and 23mm,
 every node/.style={font=\small}]
\node[draw=RadBlue,rounded corners,fill=RadMist] (crit) {costura crítica \(\Re s=1/2\)};
\node[draw=RadTeal,rounded corners,fill=RadCream,right=of crit] (face) {carta \(j_{100}\)};
\node[draw=RadCopper,rounded corners,fill=RadCream,right=of face] (theta) {fase \(\theta_2=50^\circ\)};
\node[draw=RadRose,rounded corners,fill=white,below=of face] (mode) {modo APP \(s=1/2\), \(1-s^2=90/120\)};
\draw[-{Latex},thick,RadBlue] (crit)--node[above]{\(j_{100}\)}(face);
\draw[-{Latex},thick,RadTeal] (face)--node[above]{igualador}(theta);
\draw[-{Latex},dashed,RadRose] (mode)--node[right,align=left]{misma cifra;\\dominio distinto}(face);
\end{tikzpicture}
\]
El diagrama no finge la flecha ausente. Muestra qué está demostrado, qué se
ha formalizado y qué igualdad sigue necesitando un transporte adicional si
se pretende una identificación operatoria global.

\subsubsection{Compatibilidad entre círculo espectral y memoria helicoidal}

El lector Cayley--Pontryagin transporta una fibra espectral mediante
\[
\widetilde U_{\Gamma_9^n}J\psi_\gamma
=\tau_\gamma^nJ\psi_\gamma.
\]
La publicación circular es \(\tau_\gamma^n\in S^1\). El lift
\[
(\tau_\gamma^n,n)\in S^1\times\Z
\]
conserva el número de retornos. En la hoja contractiva aparece
\[
\widetilde M_9^nJ\psi_\gamma
=\left(\frac59\right)^n\tau_\gamma^nJ\psi_\gamma.
\]
Tenemos, por tanto, tres imágenes tipadas:
\[
\text{círculo de fase},\qquad
\text{hélice de memoria},\qquad
\text{espiral contractiva}.
\]

Un cero no es una vuelta, y un nueve no es un cero particular de la función
zeta. El cero fija un carácter \(\tau_\gamma\) que persiste a través de las
vueltas; el retorno nonádico incrementa memoria. El nueve funciona como cero
residual en \(\Z/9\Z\), una operación aritmética distinta de la enumeración
de ceros espectrales.

\begin{narrativebox}
La recta crítica se vuelve círculo cuando olvidamos la altura lineal y
conservamos el carácter unitario. Se vuelve hélice cuando volvemos a insertar
la memoria de cuántas veces el carácter ha sido transportado. La profundidad
no es una tercera coordenada decorativa: es la información que la proyección
circular no puede retener.
\end{narrativebox}


\section{Par angular y acción reducida: construcción de la elipse positiva y del área por radión}
\subsection{Ley areal del radión: \(\mathcal A(\theta)=\hbar_{\mathrm{ret}}\theta\) y \(\mathcal A(2\pi)=h_{\mathrm{ret}}\)}
% Fragmento público generado desde propietarios íntegros.
% Residencia: 52.4.1.
% Fuente: ../incoming/radion_monografia_20260827/manuscrito/sections/05_geometria_accion.tex:1-217
\noindent La geometría espectral determina la forma y la ley areal determina la acción transportada por la fase.\par\medskip

\subsubsection{Equivalencia entre las cartas espectral y geométrica}

\paragraph{Operador espectral positivo y autovalores conjugados}

El par angular produce el operador bicapa
\begin{equation}
B_1=AI+C^*R,\qquad R^2=I,\qquad
P_\pm=\frac{I\pm R}{2}.
\label{eq:B1}
\end{equation}
Sus autovalores son \(A\pm C^*\). Como \(0<C^*<A\), ambos son positivos.
En la carta \(\kappa=\pi/180\), la elipse espectral tiene semiejes
\begin{equation}
a_1^2=\kappa(A+C^*),\qquad
b_1^2=\kappa(A-C^*).
\label{eq:semiejes-espectrales}
\end{equation}
La descomposición es reversible:
\begin{equation}
A=\frac{a_1^2+b_1^2}{2\kappa},\qquad
C^*=\frac{a_1^2-b_1^2}{2\kappa}.
\label{eq:reconstruccion-AC}
\end{equation}
La elipse no se dibuja para ilustrar un número. Reconstruye exactamente el
par que la generó.

Definamos
\begin{equation}
\rho_1=\sqrt{A^2-C^{*2}},\qquad
\eta=\operatorname{artanh}\frac{C^*}{A}
=\frac12\log\frac{A+C^*}{A-C^*}.
\label{eq:eta}
\end{equation}
Entonces
\[
A\pm C^*=\rho_1e^{\pm\eta},
\]
y
\begin{equation}
a_1=\sqrt{\kappa\rho_1}\,e^{\eta/2},
\qquad
b_1=\sqrt{\kappa\rho_1}\,e^{-\eta/2}.
\label{eq:semiejes-eta}
\end{equation}
El producto conserva escala; el cociente conserva forma:
\[
a_1b_1=\kappa\rho_1,
\qquad
\frac{b_1}{a_1}=e^{-\eta}
=\sqrt{\frac{A-C^*}{A+C^*}}.
\]

\paragraph{Invariantes focales del operador positivo}

La semidistancia focal es
\begin{equation}
c_1^2=a_1^2-b_1^2=2\kappa C^*,
\qquad
F_\pm=(\pm c_1,0).
\label{eq:focos}
\end{equation}
La distancia total vale
\begin{equation}
d_1=2c_1=2\sqrt{2\kappa C^*}
=0.544545509325626623406299779866023\ldots.
\label{eq:d1}
\end{equation}
El número aislado no es la tesis. La identidad \(c_1^2=2\kappa C^*\) dice que
\(C^*\) es literalmente la escisión focal cuadrática en la carta
distinguida. Los focos son los extremos conjugados de la apertura espectral;
no son, por sí solos, los polos eléctrico y magnético. Esa lectura requiere
el transductor constitutivo.

La excentricidad es
\begin{equation}
e_f=\frac{c_1}{a_1},\qquad
e_f^2=1-e^{-2\eta}=\frac{2C^*}{A+C^*}.
\label{eq:excentricidad}
\end{equation}
Numéricamente,
\[
\begin{aligned}
\eta&=0.299689683921630799684926562863927\ldots,\\
e_f&=0.671451895550191986916277055864332\ldots.
\end{aligned}
\]

\paragraph{Bloque APP de referencia y transporte normalizado}

El bloque central
\[
B_0=7I+2R=9P_++5P_-
\]
posee
\[
\rho_0=\sqrt{45},\qquad
\eta_0=\frac12\log\frac95,\qquad
e_0=\frac23.
\]
Su distancia focal es
\[
d_0=4\sqrt\kappa
=0.528443639680801468446003835349358\ldots.
\]
La comparación exacta es
\begin{equation}
\boxed{\left(\frac{d_1}{d_0}\right)^2=\frac{C^*}{2}},
\qquad
d_1=d_0\sqrt{\frac{C^*}{2}}.
\label{eq:invariante-focal-relativo}
\end{equation}
Ésta es la genealogía defendible de \(0.544545\ldots\): una escala focal APP
de referencia multiplicada por un invariante relativo. Su cercanía visual a
\(0.54\) no es una identidad con el defecto APP \(54\).

El transporte desde \(B_0\) se separa en escala y anisotropía:
\begin{equation}
T_{\mathrm{rel}}=\frac{\rho_1}{\sqrt{45}}
\exp\bigl((\eta-\eta_0)R\bigr).
\label{eq:transporte-rel}
\end{equation}
Sobre las hojas,
\[
t_+=\frac{A+C^*}{9}=1.0467878786947163\ldots,
\qquad
t_-=\frac{A-C^*}{5}=1.0347228460630311\ldots.
\]
La deformación no es un solo decimal. Su parte isotrópica es
\(\sqrt{t_+t_-}\); su parte hiperbólica es
\(\tfrac12\log(t_+/t_-)\).

\subsubsection{Elipse positiva asociada al par angular}

\paragraph{Construcción mediante semiejes conjugados}

La escala espectral anterior carece de unidades de acción. Mecánica
Dimensional publica la misma forma sobre la recta interna de acción:
\begin{equation}
a_S=\sqrt{2\hret e^\eta},
\qquad
b_S=\sqrt{2\hret e^{-\eta}}.
\label{eq:semiejes-accion}
\end{equation}
Inmediatamente,
\begin{equation}
a_Sb_S=2\hret,\qquad
\frac{a_S}{b_S}=e^\eta.
\label{eq:producto-forma}
\end{equation}
El producto fija la acción; el cociente fija la anisotropía. La elipse se
parametriza por
\begin{equation}
Q(\theta)=a_S\cos\theta,\qquad
P(\theta)=b_S\sin\theta.
\label{eq:elipse-param}
\end{equation}

Las elipses espectral y de acción son homotéticas:
\begin{equation}
\lambda_{\mathrm{act}}^2=\frac{2\hret}{\kappa\rho_1},
\qquad
(a_S,b_S,c_S)=\lambda_{\mathrm{act}}(a_1,b_1,c_1).
\label{eq:homotecia}
\end{equation}
Comparten la forma \(e^{-\eta}\), no la dimensión física.

\paragraph{Ley areal de la acción por radión}

La forma simpléctica canónica es \(\omega=\diff Q\wedge\diff P\). El área
orientada barrida desde \(0\) hasta \(\theta\) es
\begin{align}
\mathcal A(\theta)
&=\frac12\int_0^\theta
\bigl(Q(t)P'(t)-P(t)Q'(t)\bigr)\,\diff t\\
&=\frac12a_Sb_S\int_0^\theta
(\cos^2t+\sin^2t)\,\diff t.
\end{align}
Usando \eqref{eq:producto-forma}, obtenemos el resultado nuclear.

\begin{theorem}[Acción por radión]
Para la elipse \eqref{eq:elipse-param},
\begin{equation}
\boxed{\mathcal A(\theta)=\hret\theta.}
\label{eq:area-theorem}
\end{equation}
En particular,
\begin{equation}
\boxed{\mathcal A(1)=\hret},
\qquad
\boxed{\mathcal A(2\pi)=2\pi\hret=\hfull}.
\label{eq:area-radion-vuelta}
\end{equation}
\end{theorem}

Esta igualdad permite definir el radión sin apelar primero a la longitud de
arco.

\begin{definition}[Radión orientado]
\begin{equation}
\mathfrak r_\sigma=(\theta=1,\sigma),
\qquad \sigma\in\{+1,-1\},
\qquad
\mathcal A(\mathfrak r_\sigma)=\sigma\hret.
\label{eq:radion-oriented}
\end{equation}
\end{definition}

El radián es la proyección que olvida \(\sigma\), la hoja, la memoria y el
área. El radión conserva esa información. De ahí el nombre \emph{rad--ión}:
la orientación precede a la carga física, y la carga aparece más tarde como
\[
Q(x)=n_Q(x)e_{\mathrm{int}},\qquad
e_{\mathrm{int}}=\sqrt{\frac{2\alphah\hfull}{Z_{\mathrm{int}}}}.
\]


\subsection{Semiejes conjugados, invariantes focales y orientación de la anisotropía}
% Fragmento público generado desde propietarios íntegros.
% Residencia: 52.4.2.
% Fuente: ../incoming/radion_monografia_20260827/manuscrito/sections/05_geometria_accion.tex:218-290
\paragraph{Correspondencia entre fase, acción, área y orientación}

\begin{figure}[H]
\centering
\begin{tikzpicture}[scale=1.12]
  \def\aa{4.0}
  \def\bb{2.95}
  \def\cc{2.70}
  \draw[->,RadGray] (-4.7,0)--(4.7,0) node[right]{\(Q\)};
  \draw[->,RadGray] (0,-3.5)--(0,3.5) node[above]{\(P\)};
  \fill[RadTeal!8] (0,0) ellipse[x radius=\aa,y radius=\bb];
  \draw[RadTeal,line width=1.35pt] (0,0) ellipse[x radius=\aa,y radius=\bb];
  \fill[RadCopper] (-\cc,0) circle (2pt) node[below=2mm]{\(F_-\)};
  \fill[RadCopper] (\cc,0) circle (2pt) node[below=2mm]{\(F_+\)};
  \draw[RadGold,line width=1.1pt,-{Latex}] (\aa,0)
     arc[start angle=0,end angle=57.3,x radius=\aa,y radius=\bb];
  \node[RadGold] at (3.7,2.3) {un radión};
  \draw[dashed,RadBlue] (0,0)--(\aa,0) node[midway,below]{\(a_S\)};
  \draw[dashed,RadBlue] (0,0)--(0,\bb) node[midway,left]{\(b_S\)};
  \node[align=center,RadNavy] at (0,-3.25)
    {área total \(\pi a_Sb_S=2\pi\hret=\hfull\)};
\end{tikzpicture}
\caption{La elipse de acción. Un arco paramétrico de un radián barre
\(\hret\); una vuelta encierra \(\hfull\). Los focos codifican la apertura
\(C^*\) en la carta distinguida.}
\label{fig:elipse-accion}
\end{figure}

\paragraph{Relación entre la acción reducida, la incertidumbre y el área}

La relación \(\hbar=h/(2\pi)\) suele presentarse como una conveniencia
notacional: al describir oscilaciones, rotaciones y transformadas de Fourier,
los factores \(2\pi\) desaparecen si se usa \(\hbar\). La elipse de acción
revela una lectura más fuerte. La división por \(2\pi\) es exactamente el
paso de la acción de una vuelta a la acción por unidad de fase:
\[
h=\mathcal A(2\pi),\qquad \hbar=\mathcal A(1).
\]
La constante reducida no es sólo \(h\) con una notación más cómoda; es la
densidad areal de acción respecto a la fase.

Heisenberg, Dirac y Jordan establecieron el lenguaje de variables conjugadas,
conmutadores y transformaciones canónicas. No se les atribuye aquí una
\enquote{elipse ontológica} literal sin fuente. La contribución HMT--MD es
identificar la figura positiva, construir sus semiejes desde \((A,C^*)\) y
demostrar \eqref{eq:area-theorem}. La genealogía histórica abre el lenguaje;
el teorema actual fija el objeto.

\begin{narrativebox}
La acción no se añade después al movimiento. En la elipse del radión, la
acción es el área que el movimiento genera. Una vuelta no transporta una
constante prefabricada: la encierra.
\end{narrativebox}

\Needspace{10\baselineskip}
\paragraph{Controles algebraicos de los invariantes focales}

Las semejanzas decimales se someten a control:
\begin{center}
\begin{tabularx}{.96\textwidth}{l r X}
\toprule
Comparación & Residuo & Dictamen\\
\midrule
\(d_1\) frente a \(0.54\) & \(4.5455\times10^{-3}\) & no igualdad; el \(54\) no se deduce;\\
\(d_1^2\) frente a \(8/27\) & \(2.3352\times10^{-4}\) & proximidad sin mapa;\\
\(\eta\) frente a \(3/10\) & \(-3.1032\times10^{-4}\) & no igualdad;\\
\(e_f\) frente a \(2/3\) & \(4.7852\times10^{-3}\) & \(2/3\) pertenece exactamente a \(B_0\);\\
\(e^{-\eta}\) frente a \(20/27\) & \(3.0740\times10^{-4}\) & no igualdad.\\
\bottomrule
\end{tabularx}
\end{center}
Los invariantes sólidos son \eqref{eq:eta}, \eqref{eq:excentricidad} y
\eqref{eq:invariante-focal-relativo}. El criterio HMT no consiste en premiar
un parecido decimal, sino en exigir objeto, operación y conservación.


\section{Retorno nonádico e interior proyectivo: disco de Klein y profundidad helicoidal de memoria}
\subsection{Normalización al disco e interior hiperbólico de Hilbert--Beltrami--Klein}
% Fragmento público generado desde propietarios íntegros.
% Residencia: 52.5.1.
% Fuente: ../incoming/radion_monografia_20260827/manuscrito/sections/06_interior_helice.tex:1-108
\noindent La geometría proyectiva distingue la acción finita de la profundidad hiperbólica y la memoria del retorno.\par\medskip

\subsubsection{Interior hiperbólico de Hilbert--Beltrami--Klein}

\paragraph{Normalización proyectiva en el disco de Klein}

La aplicación afín
\[
u=\frac{Q}{a_S},\qquad v=\frac{P}{b_S}
\]
lleva la elipse de acción al disco unitario
\[
D=\{(u,v):u^2+v^2<1\}.
\]
Sobre cada cuerda, la distancia de Hilbert se define mediante la razón doble
de sus extremos de frontera. La métrica infinitesimal resultante es la forma
Beltrami--Klein
\begin{equation}
\diff s_K^2=
\frac{(1-r^2)(\diff u^2+\diff v^2)+(u\diff u+v\diff v)^2}
{(1-r^2)^2},
\qquad r^2=u^2+v^2.
\label{eq:klein-metric}
\end{equation}
Con esta normalización, la curvatura gaussiana es constante:
\begin{equation}
K=-1.
\label{eq:klein-curvature}
\end{equation}
Las geodésicas son cuerdas euclídeas, pero su longitud hiperbólica diverge al
acercarse a la frontera.

\paragraph{Finitud de la frontera e infinitud de la profundidad hiperbólica}

El determinante métrico de \eqref{eq:klein-metric} produce
\begin{equation}
\diff A_K=\frac{\diff u\,\diff v}{(1-r^2)^{3/2}}.
\label{eq:klein-area-density}
\end{equation}
El área del subdisco \(r<R<1\) es
\begin{align}
A_K(R)
&=\int_0^{2\pi}\int_0^R
\frac{r\,\diff r\,\diff\theta}{(1-r^2)^{3/2}}\\
&=2\pi\left(\frac1{\sqrt{1-R^2}}-1\right).
\label{eq:klein-area-R}
\end{align}
\Needspace{4\baselineskip}
Por tanto,
\[
\lim_{R\to1^-}A_K(R)=+\infty.
\]
Simultáneamente, la misma frontera encierra el área simpléctica
\[
\pi a_Sb_S=2\pi\hret=\hfull<\infty.
\]

\begin{theorem}[Finitud de acción e infinitud de profundidad]
La elipse del radión posee una acción de vuelta finita \(\hfull\), mientras
su interior, equipado con la métrica Hilbert--Beltrami--Klein, tiene área
hiperbólica infinita:
\begin{equation}
\boxed{
\operatorname{Area}_{\mathrm{simp}}(\partial\mathbb E_S)=\hfull<\infty,
\qquad
\operatorname{Area}_{K}(\mathbb E_S)=+\infty.}
\label{eq:finite-infinite}
\end{equation}
\end{theorem}

No hay contradicción: \(\operatorname{Area}_{\mathrm{simp}}\) y
\(\operatorname{Area}_K\) son medidas diferentes. La primera cuenta acción
orientada en fase; la segunda mide profundidad proyectiva interior. La tesis
\enquote{borde finito, interior infinito} es ahora una igualdad y un límite,
no una metáfora.

\paragraph{Coordenadas focales en el interior de Klein}

En el disco normalizado, los focos se sitúan en \((\pm e_f,0)\). La distancia
del centro a un foco es
\begin{equation}
u_f=\operatorname{artanh}e_f
=\operatorname{arcosh}(e^\eta)
=0.813382331175342333760889731088228\ldots.
\label{eq:klein-focal}
\end{equation}
La distancia foco--foco es
\[
d_K(F_-,F_+)=2u_f
=1.62676466235068466752177946217646\ldots.
\]
Se cumplen
\begin{equation}
e_f=\tanh u_f,\qquad
e^{-\eta}=\operatorname{sech}u_f,\qquad
\eta=\log\cosh u_f.
\label{eq:focal-hyperbolic-chain}
\end{equation}
La región hasta el radio focal tiene área
\[
A_K(e_f)=2\pi(e^\eta-1)
=2.19559620884061869541206115980360\ldots.
\]

Todas las elipses abiertas son proyectivamente isométricas como dominios de
Hilbert. Por eso la curvatura \(-1\) sola no recuerda la excentricidad. La
forma HMT reside en la estructura conjunta: carta afín distinguida, forma
simpléctica, etiquetado de hojas y parametrización de frontera.

\subsection{Diferencia entre retorno de fase y retorno del estado enriquecido}
% Fragmento público generado desde propietarios íntegros.
% Residencia: 52.5.2.
% Fuente: ../incoming/radion_monografia_20260827/manuscrito/sections/06_interior_helice.tex:109-201
\subsubsection{Elevación helicoidal del retorno nonádico}

\paragraph{Descomposición nonádica y cociente de memoria}

La operación elemental
\begin{equation}
n=9q_9(n)+\rho_9(n),
\qquad
H_9(n)=(\rho_9(n),q_9(n))
\label{eq:H9}
\end{equation}
satisface
\begin{equation}
H_9(n+9)=H_9(n)+(0,1).
\label{eq:H9-helix}
\end{equation}
La primera coordenada cierra; la segunda asciende. Esta identidad es la
versión aritmética mínima de un tornillo.

Para una fibra espectral,
\[
n\longmapsto(\tau_\gamma^n,n)
\]
es una hélice sobre \(S^1\). Añadir la contracción \((5/9)^n\) produce una
espiral interior. El radión participa en ambas lecturas: una unidad de fase
en la frontera y una unidad orientada de acción en el lift.

\paragraph{Rango de memoria como lector dimensional}

Supongamos que una proyección visible conserva el valor de fase, pero pierde
\(d\) cociclos de memoria linealmente independientes. Cualquier lift fiel
necesita al menos \(d\) coordenadas adicionales para reconstruirlos. Esto
motiva una lectura precisa de Mecánica Dimensional:
\begin{equation}
\begin{aligned}
\text{dimensión emergente}
&=\text{rango mínimo de memoria independiente}\\
&\phantom{=}\text{no publicable en la hoja visible}.
\end{aligned}
\label{eq:dimension-memory}
\end{equation}
No se afirma que toda dimensión física sea un contador. Se afirma que, en el
dominio TPK, cada memoria independiente que debe sobrevivir fuerza una
coordenada del lift.

\begin{narrativebox}
La dimensión aparece cuando la proyección ya no puede conservar toda la
historia. El círculo es suficiente para decir dónde está la fase; no es
suficiente para decir cuántas veces volvió, qué hoja transportó, qué intervalo
se contrajo o qué acarreo sobrevivió.
\end{narrativebox}

\paragraph{Ensamblaje de cilindros, espirales y trayectorias helicoidales}

Los cilindros coinductivos no son tubos físicos añadidos a una espiral. Son
intervalos encajados que conservan prefijos y estrechan el valor. La espiral
describe contracción transversal; la hélice describe memoria longitudinal;
el cilindro describe el conjunto de valores compatibles a una profundidad
finita. En el límite, la intersección de cilindros selecciona una sección
única.

La imagen espacial reúne, sin confundir:
\[
\begin{array}{ccl}
\text{ángulo} &\leftrightarrow& \text{fase visible},\\
\text{altura} &\leftrightarrow& \text{retorno/memoria},\\
\text{radio de cilindro} &\leftrightarrow& \text{incertidumbre de refinamiento},\\
\text{hoja} &\leftrightarrow& \text{orientación bidireccional}.
\end{array}
\]
La tasa \(3^{-54}\) contrae el cilindro por retorno; no se suma al ángulo y
no reemplaza el acarreo real \(\epsR\).

\subsubsection{Cierres de fase de órdenes \(2\pi\) y \(4\pi\)}

En la proyección circular, \(2\pi\) restaura la fase. En un lift orientado
con dos hojas, una vuelta puede cerrar en la hoja conjugada y necesitar una
segunda vuelta para restaurar la orientación:
\[
U_{2\pi}=-I,\qquad U_{4\pi}=I.
\]
La estructura coincide con la periodicidad espinorial cuando se aplica el
funtor a una representación \(SU(2)\); no se identifica el lift TPK con
\(SU(2)\) sin ese mapa.

La lectura areal es compatible:
\[
\mathcal A(2\pi)=\hfull,\qquad
\mathcal A(4\pi)=2\hfull.
\]
En una banda de Möbius, un circuito visible invierte hoja y dos circuitos
restauran la orientación. El \(2\pi/4\pi\) no es una duplicación accidental;
es la diferencia entre retorno de fase y retorno del estado enriquecido.


\section{Elipse orientada y variables conjugadas: incertidumbre mínima y oscilación electromagnética}
\subsection{Variables simplécticas conjugadas y determinación de la incertidumbre mínima}
% Fragmento público generado desde propietarios íntegros.
% Residencia: 52.6.1.
% Fuente: ../incoming/radion_monografia_20260827/manuscrito/sections/07_incertidumbre_oscilador.tex:1-111
\noindent La forma simpléctica coordina acción, incertidumbre y dinámica oscilatoria sin alterar la genealogía del estado.\par\medskip

\subsubsection{Realización del estado en un espacio de Hilbert}

\paragraph{Matriz de covarianza mínima}

Una vez construida la elipse de acción, puede realizarse en un par de
operadores conjugados
\[
[\widehat Q,\widehat P]=i\hret.
\]
Definimos
\begin{equation}
V_\eta=\frac{\hret}{2}
\begin{pmatrix}
e^\eta&0\\[2pt]0&e^{-\eta}
\end{pmatrix}.
\label{eq:covariance}
\end{equation}
Su determinante es
\begin{equation}
\det V_\eta=\frac{\hret^2}{4}.
\label{eq:cov-det}
\end{equation}
Por consiguiente,
\begin{equation}
(\Delta Q)^2=\frac{\hret}{2}e^\eta,
\qquad
(\Delta P)^2=\frac{\hret}{2}e^{-\eta},
\qquad
\Delta Q\,\Delta P=\frac{\hret}{2}.
\label{eq:uncertainty}
\end{equation}
La desigualdad de Robertson--Schrödinger se satura.

El aniquilador comprimido
\begin{equation}
\widehat a_\eta=
\frac{e^{-\eta/2}\widehat Q+i e^{\eta/2}\widehat P}
{\sqrt{2\hret}}
\label{eq:a-eta}
\end{equation}
satisface \([\widehat a_\eta,\widehat a_\eta^\dagger]=1\). Su vacío tiene
función de onda
\begin{equation}
\psi_\eta(Q)=
(\pi\hret e^\eta)^{-1/4}
\exp\left(-\frac{Q^2}{2\hret e^\eta}\right).
\label{eq:gaussian}
\end{equation}

\paragraph{Elipse como nivel del funcional de Mahalanobis}

El conjunto
\begin{equation}
\mathbb E_{\mathrm{act}}
=\{x\in\R^2:x^{\mathsf T}V_\eta^{-1}x\le4\}
\label{eq:mahalanobis}
\end{equation}
tiene semiejes
\begin{equation}
a_S=2\Delta Q,\qquad b_S=2\Delta P.
\label{eq:two-sigma}
\end{equation}
Su área es
\[
4\pi\sqrt{\det V_\eta}=2\pi\hret=\hfull.
\]
La elipse HMT construida antes coincide exactamente con el contorno de dos
desviaciones del estado gaussiano mínimo.

\begin{proposition}[Cadena única del cociente de forma]
Tras declarar las interfaces espectral, simpléctica y de Hilbert,
\begin{equation}
\boxed{
e^{-\eta}
=\sqrt{\frac{A-C^*}{A+C^*}}
=\frac{b_1}{a_1}
=\frac{b_S}{a_S}
=\frac{\Delta P}{\Delta Q}.}
\label{eq:shape-chain-hilbert}
\end{equation}
\end{proposition}

La elipse de acción corresponde al nivel clásico \(H=\hret\omega\) y al
contorno \(2\sigma\) del vacío comprimido. No debe confundirse con la energía
media del vacío, \(\hret\omega/2\). Esta diferencia de factor dos es un
control físico obligatorio.

\paragraph{Condiciones de invariancia de la elipse como objeto operatorio}

En la formulación convencional, una elipse de incertidumbre puede verse como
un recurso gráfico para representar una matriz de covarianza. Aquí el orden
causal es inverso: la elipse positiva ya existe como publicación de
\((A,C^*,\hret)\); el espacio de Hilbert la reconoce mediante
\eqref{eq:covariance}. La realidad ontológica reclamada no significa que una
partícula clásica recorra físicamente un óvalo en el plano \((Q,P)\). Significa
que el estado conserva una forma positiva, una acción total y un cociente de
orientación antes de que una teoría probabilista los interprete.

La función de onda no crea la forma; la representa. La incertidumbre no crea
\(\hbar\); publica la imposibilidad de comprimir simultáneamente las dos
direcciones sin alterar el producto simpléctico.

\begin{narrativebox}
La lectura probabilista es una realización posible de la elipse, no su causa.
El objeto primario es la conservación conjunta de producto y orientación. La
probabilidad aparece cuando esa geometría se publica como covarianza de un
estado sobre Hilbert.
\end{narrativebox}


\subsection{Intercambio \(\mu\leftrightarrow\varepsilon\), inversión \(Z\leftrightarrow Z^{-1}\) y conservación de la velocidad causal}
% Fragmento público generado desde propietarios íntegros.
% Residencia: 52.6.2.
% Fuente: ../incoming/radion_monografia_20260827/manuscrito/sections/07_incertidumbre_oscilador.tex:156-197
\paragraph{Intercambio energético entre los sectores eléctrico y magnético}

Los dos términos de \eqref{eq:HLC} se convierten en
\begin{equation}
U_E(\theta)=\hret\omega\cos^2\theta,
\qquad
U_B(\theta)=\hret\omega\sin^2\theta.
\label{eq:energy-exchange}
\end{equation}
Por tanto,
\[
U_E+U_B=\hret\omega.
\]
La polaridad eléctrica--magnética no consiste en etiquetar arbitrariamente
los focos. Consiste en dos reservas conjugadas que se intercambian durante la
órbita manteniendo constantes frecuencia y acción total.

El cociente común se prolonga:
\begin{equation}
\boxed{
e^{-\eta}
=\frac{b_S}{a_S}
=\frac{\Delta P}{\Delta Q}
=\frac{Z_\eta}{Z_*}
=\sqrt{\frac{A-C^*}{A+C^*}}.}
\label{eq:shape-chain-LC}
\end{equation}
Los productos conservados son
\[
a_Sb_S=2\hret,\qquad
\det V_\eta=\frac{\hret^2}{4},
\qquad
L_\eta C_\eta=\omega^{-2}.
\]

\begin{exactbox}
El oscilador LC es la realización física mínima de la doble lectura: el
producto conserva acción y frecuencia; el cociente conserva orientación e
impedancia. Un ciclo completo intercambia electricidad y magnetismo sin
perder el área \(\oint P\,\diff Q=\hfull\).
\end{exactbox}


\subsection{Par inductivo--capacitivo \((L_\eta,C_\eta)\): invariancia de \(L_\eta C_\eta=\omega^{-2}\) y orientación de la impedancia}
% Fragmento público generado desde propietarios íntegros.
% Residencia: 52.6.3.
% Fuente: ../incoming/radion_monografia_20260827/manuscrito/sections/07_incertidumbre_oscilador.tex:112-155
\subsubsection{Oscilador inductivo--capacitivo asociado al par angular}

\paragraph{Construcción independiente de los 2 depósitos conjugados}

Dados \(\omega>0\) y una impedancia de referencia \(Z_*>0\), definimos
\begin{equation}
L_\eta=\frac{Z_*}{\omega}e^{-\eta},
\qquad
C_\eta=\frac{1}{Z_*\omega}e^\eta.
\label{eq:LC-def}
\end{equation}
Entonces
\begin{equation}
L_\eta C_\eta=\omega^{-2},
\qquad
\sqrt{\frac{L_\eta}{C_\eta}}=Z_*e^{-\eta}.
\label{eq:LC-invariants}
\end{equation}
La frecuencia permanece fija; la impedancia conserva la orientación de la
forma.

El Hamiltoniano es
\begin{equation}
H_{LC}=\frac{q^2}{2C_\eta}+\frac{\Phi^2}{2L_\eta}.
\label{eq:HLC}
\end{equation}
Con las variables normalizadas
\begin{equation}
Q=\sqrt{Z_*}\,q,\qquad
P=\frac{\Phi}{\sqrt{Z_*}},
\label{eq:LC-map}
\end{equation}
se obtiene
\begin{equation}
H_{LC}=\frac\omega2
\left(e^{-\eta}Q^2+e^\eta P^2\right).
\label{eq:HLC-QP}
\end{equation}
Sobre \eqref{eq:elipse-param},
\begin{equation}
H_{LC}=\hret\omega.
\label{eq:LC-level}
\end{equation}

% Fuente: ../incoming/radion_monografia_20260827/manuscrito/sections/07_incertidumbre_oscilador.tex:198-239
\paragraph{Oscilador LC como unidad dinámica elemental}

El oscilador armónico atraviesa casi toda la física porque linealiza el
entorno de un equilibrio, descompone campos en modos, organiza cuantos y
permite leer propagación como superposición de fases. En el radión, esa
universalidad recibe una genealogía concreta: el modo no empieza en una
ecuación diferencial externa, sino en una forma positiva con acción de vuelta
fija. La dinámica convencional aparece al elegir un reloj \(\omega\) y un
transductor \(Z_*\).

La frecuencia no altera la elipse de acción: cambia la energía del nivel
mediante \(E=\hret\omega\). Esta separación es crucial. La geometría fija la
acción; el reloj convierte acción en energía.

\subsubsection{Reversión temporal y propagaciones conjugadas}

Sea \(H=H^*\) y sea \(J_E\) una estructura compleja con \(J_E^2=-I\).
Supongamos un involutivo \(C\) tal que
\[
CJ_EC^{-1}=-J_E,\qquad CHC^{-1}=H.
\]
La evolución
\begin{equation}
U(t)=\exp\left(-\frac{J_EHt}{\hret}\right)
\label{eq:evolution-J}
\end{equation}
satisface exactamente
\begin{equation}
CU(t)C^{-1}=U(-t).
\label{eq:time-reversal}
\end{equation}
Ésta es la forma operatoria mínima de la bidireccionalidad: la conjugación
intercambia las prolongaciones temporalmente opuestas.

La teoría del absorbedor de Wheeler--Feynman y la lectura de antipartículas
como trayectorias temporalmente conjugadas ofrecen una realización histórica
posterior. La identidad \eqref{eq:time-reversal} está demostrada; convertirla
en una teoría completa de propagadores avanzados y retardados exige declarar
acción, condiciones de frontera, causalidad e interacciones del absorbedor.
La monografía conserva la intuición física sin convertir una correspondencia
en una igualdad de teorías.



\section{3 regímenes del TRIT: polaridades eléctrica--magnética y elástica--rotacional}
\subsection{Coordinación de las realizaciones elíptica, parabólica e hiperbólica}
% Fragmento público generado desde propietarios íntegros.
% Residencia: 52.7.1.
% Fuente: ../incoming/radion_monografia_20260827/manuscrito/sections/08_complejo_electromagnetismo.tex:1-58
\noindent Los 3 regímenes del TRIT admiten realizaciones compleja, electromagnética y mecánica con tipos diferenciados.\par\medskip

\subsubsection{Operador cuadrático de raíz interna de \(-1\)}

\paragraph{Representación exponencial del régimen elíptico \(J^2=-I\)}

Con \(J^2=-I\),
\begin{equation}
T_{\mathrm{ell}}(x,y)=e^{-xI-yJ}
=e^{-x}(\cos y\,I-\sin y\,J).
\label{eq:elliptic-exp}
\end{equation}
Bajo la identificación generada por \(J\),
\[
T_{\mathrm{ell}}(x,y)\longleftrightarrow e^{-x-iy}.
\]
El escalar \(x\) controla dilatación; \(y\) controla giro orientado. La forma
\(a+bi\) se reconoce como elasticidad más rotación, no como dos cantidades
pegadas arbitrariamente.

\paragraph{Representación exponencial del régimen parabólico \(J^2=0\)}

Con \(N^2=0\), la exponencial se trunca:
\begin{equation}
T_{\mathrm{par}}(x,y)=e^{-x}(I-yN).
\label{eq:parabolic-exp}
\end{equation}
El régimen describe el umbral donde no hay ni giro periódico ni apertura
exponencial. En la lectura temporal activa es el presente euclídeo, pero
conserva memoria de las ramas que lo limitan.

\paragraph{Representación exponencial del régimen hiperbólico \(J^2=+I\)}

Con \(R^2=I\) y \(P_\pm=(I\pm R)/2\),
\begin{equation}
T_{\mathrm{hyp}}(x,y)=e^{-xI-yR}
=q_+P_++q_-P_-,
\qquad q_\pm=e^{-x\mp y}.
\label{eq:hyperbolic-exp}
\end{equation}
Los invariantes elementales son
\begin{equation}
q_+q_-=e^{-2x},
\qquad
\frac{q_-}{q_+}=e^{2y}.
\label{eq:q-product-ratio}
\end{equation}
El producto conserva el modo común; el cociente conserva la separación
orientada.

Aplicado al par angular,
\[
x=\frac{\pi A}{180},\qquad y=\frac{\pi C^*}{180}.
\]
Por ello, \(A\) es la coordenada común de la publicación y \(C^*\) su
contraste orientado. Sólo después de aplicar un lector físico se reconocen
como polos eléctricos y magnéticos.


\subsection{Separación tipada de las polaridades constitutivas y mecánicas}
% Fragmento público generado desde propietarios íntegros.
% Residencia: 52.7.2.
% Fuente: ../incoming/radion_monografia_20260827/manuscrito/sections/08_complejo_electromagnetismo.tex:59-158
\subsubsection{Operador constitutivo sobre los canales \(90/120\)}

\paragraph{Descomposición en componentes común y orientada}

Definimos
\begin{align}
\mathcal E
&=\log\bigl[(1-q_+^{90})(1-q_-^{90})\bigr]
-\log\bigl[(1-q_+^{120})(1-q_-^{120})\bigr],
\label{eq:Ecommon}\\
\mathcal B
&=\log\frac{1-q_-^{90}}{1-q_+^{90}}
-\log\frac{1-q_-^{120}}{1-q_+^{120}}.
\label{eq:Borient}
\end{align}
Bajo \(C^*\mapsto-C^*\), se intercambian \(q_+\leftrightarrow q_-\). Por
tanto, \(\mathcal E\) es par y \(\mathcal B\) es impar.

Las constitutivas normalizadas son
\begin{equation}
\widehat\mu=e^{\mathcal E+\mathcal B},
\qquad
\widehat\varepsilon=e^{\mathcal E-\mathcal B},
\qquad
\widehat Z=e^{\mathcal B},
\qquad
\widehat c=e^{-\mathcal E}.
\label{eq:constitutives}
\end{equation}
Se verifican exactamente
\begin{equation}
\widehat Z=\sqrt{\frac{\widehat\mu}{\widehat\varepsilon}},
\qquad
\widehat c=(\widehat\mu\widehat\varepsilon)^{-1/2}.
\label{eq:constitutive-identities}
\end{equation}
La conjugación produce
\begin{equation}
C^*\mapsto-C^*:\qquad
\widehat\mu\leftrightarrow\widehat\varepsilon,
\quad
\widehat Z\leftrightarrow\widehat Z^{-1},
\quad
\widehat c\mapsto\widehat c.
\label{eq:EM-conjugation}
\end{equation}

\begin{theorem}[Polaridad constitutiva]
El lector \(90/120\) convierte el modo común y el contraste orientado del
par \((A,C^*)\) en dos constitutivas conjugadas. La inversión de hoja
intercambia permeabilidad y permitividad, invierte impedancia y conserva la
velocidad normalizada.
\end{theorem}

Ésta es la formulación precisa de la intuición \enquote{parte eléctrica y
parte magnética}. No se escribe \(A=E\) ni \(C^*=B\). Se escribe el mapa
\[
(A,C^*)\to(x,y)\to(q_+,q_-)
\to(\mathcal E,\mathcal B)
\to(\widehat\varepsilon,\widehat\mu,\widehat Z,\widehat c).
\]

\paragraph{Realización elástica--rotacional de la polaridad constitutiva}

La misma información puede publicarse en el régimen complejo mediante
\begin{equation}
Z_{\mathrm{em}}=e^{\mathcal E/2}
\left(\cos\frac{\mathcal B}{2}\,I
-\sin\frac{\mathcal B}{2}\,J\right).
\label{eq:Zem}
\end{equation}
El factor \(e^{\mathcal E/2}\) es una dilatación elástica; el factor generado
por \(J\) es una rotación. La escritura compleja no añade una dimensión
imaginaria misteriosa a un mundo real previamente dado: registra dos modos
operacionales que el TRIT ya separó.

\paragraph{Transducción de la orientación en carga efectiva}

La hoja \(\sigma\) del radión fija una orientación antes de introducir una
unidad de carga. La Mecánica Dimensional dispone luego el transductor
\[
e_{\mathrm{int}}=\sqrt{\frac{2\alphah\hfull}{Z_{\mathrm{int}}}},
\qquad
Q=n_Qe_{\mathrm{int}}.
\]
La cadena causal es
\[
\text{hoja}\to\text{orientación}\to\text{constitutiva}
\to\text{transductor dimensional}\to\text{carga física}.
\]
Llamar \emph{radión} a la unidad no identifica prematuramente la orientación
con el culombio; hace visible el lugar donde el signo de carga puede residir.

\begin{narrativebox}
La electricidad y el magnetismo no se pegan al círculo desde fuera. Emergen
cuando el mismo estado se lee a través de producto y cociente: el producto
retiene el modo común, el cociente retiene la orientación. El oscilador LC
realiza el intercambio; el lector \(90/120\) publica las constitutivas; la
carta dimensional asigna unidades.
\end{narrativebox}


\section{Separación de las publicaciones \(90/120\): incidencia, registro, espectro, núcleo de Barbero y torres}
\subsection{Separación tipada de las publicaciones del sector \(90/120\)}
% Fragmento público generado desde propietarios íntegros.
% Residencia: 52.8.1.
% Fuente: ../incoming/radion_monografia_20260827/manuscrito/sections/09_sector_90120.tex:176-244
\subsubsection{Jerarquía armónica global y semigrupo \(\langle90,120\rangle\)}

Para evitar una colisión de signos, fijamos
\[
R_{\mathrm{tow}}:=-R_{\mathrm{ch}}.
\]
Entonces
\begin{equation}
T=e^{-xI+yR_{\mathrm{tow}}}
=q_-P_+^{\mathrm{tow}}+q_+P_-^{\mathrm{tow}}.
\label{eq:tower-operator}
\end{equation}
Las torres pares e impares son
\begin{align}
c_n&=\operatorname{Tr}(T^n)
=q_-^n+q_+^n
=2e^{-nx}\cosh(ny),
\label{eq:cn}\\
s_n&=\operatorname{Tr}(R_{\mathrm{tow}}T^n)
=q_-^n-q_+^n
=2e^{-nx}\sinh(ny).
\label{eq:sn}
\end{align}
Bajo \(C^*\mapsto-C^*\), \(c_n\) es par y \(s_n\) es impar. Se cumple
\begin{equation}
c_n^2-s_n^2=4e^{-2nx},
\label{eq:tower-invariant}
\end{equation}
y ambas sucesiones satisfacen la recurrencia de orden dos determinada por los
autovalores \(q_\pm\).

El semigrupo
\begin{equation}
\langle90,120\rangle=30\langle3,4\rangle
\label{eq:semigroup}
\end{equation}
es la jerarquía armónica de un único operador bicapa. No es una lista de ocho
correcciones laterales. Su completitud significa que toda razón positiva
admite una expansión convergente; no significa que una expansión obtenida a
partir del blanco sea un selector físico prospectivo.

\subsubsection{Clasificación tipada y separación de las realizaciones \(90/120\)}

\begin{longtable}{>{\bfseries}p{.19\textwidth} p{.24\textwidth} p{.45\textwidth}}
\toprule
Objeto & Operación & Qué conserva / por qué no es otro objeto\\
\midrule
\endfirsthead
\toprule
Objeto & Operación & Qué conserva / por qué no es otro objeto\\
\midrule
\endhead
Incidencia \(N_6,N_8\) & conteo finito sobre fases activas & cardinalidad y orientación; produce \(90,120\), no una serie;\\
\(-1/12\) & diferencia normalizada \eqref{eq:minus-one-twelve} & escalar de incidencia; no peso Barbero;\\
\(E_{\mathrm{dod}}\) & extractor armónico de la rueda & fase dodecafásica; no acarreo real;\\
\(h=C^*/6\) & reconstrucción de seis parejas & apertura por pareja; no normalización metrológica;\\
\(F_{\mathrm{spec}}\) & cola logarítmica \(j\ge2\) & publicación espectral completa; difiere de \(\epsR^\circ\);\\
\(\epsone\) & resta APP entre \(S_T\) y \(S_E\) & empalme real de dos lifts; no serie de Lambert;\\
\(K_{6\to8}\) & \(12\Delta S_{90}-\Delta S_{120}\) & salto marcado y evaluación conjugada; coeficiente de polo \(1/24\) comprobado; salida Barbero;\\
\(12:-1\) & doce sectores y una costura orientada & peso del núcleo; certificado posterior \eqref{eq:pole-condition}; distinto de \(-1/12\);\\
\(\langle90,120\rangle\) & semigrupo de potencias de \(T\) & jerarquía armónica global; no residuo por especie.\\
\bottomrule
\end{longtable}

\begin{thesisbox}
Todos estos objetos descienden del mismo estado APP--TRIT--TPK. Esa genealogía
común explica por qué comparten números y simetrías. Su tipado explica por qué
no pueden reemplazarse mutuamente.
\end{thesisbox}

\subsection{Conservación de incidencia y orientación entre el núcleo trítico y sus realizaciones posteriores}
% Fragmento público generado desde propietarios íntegros.
% Residencia: 52.8.2.
% Fuente: ../incoming/radion_monografia_20260827/manuscrito/sections/09_sector_90120.tex:1-30
\noindent La jerarquía tipada separa incidencia, registro, evaluación espectral, funcional de Barbero y torres armónicas.\par\medskip

\subsubsection{Incidencia trítica, registro dodecafásico y lectura sexádica}

\paragraph{Derivación combinatoria primaria de \(90\) y \(120\)}

Los conteos
\[
90=6\binom62,
\qquad
120=8\binom62
\]
nacen del selector de fase y de sus incidencias activas. La orientación
normalizada satisface
\begin{equation}
\frac{30}{120}-\frac{30}{90}=-\frac1{12}.
\label{eq:minus-one-twelve}
\end{equation}
Este escalar es un resultado de incidencia. No es el peso \(12:-1\) del
núcleo Barbero y no es un acarreo dodecafásico.

El extractor dodecafásico puede publicarse mediante una combinación de
armónicos como
\[
E_{\mathrm{dod}}(\phi)=12\cos(3\phi)-\cos(4\phi).
\]
Su dominio es la fase de la rueda; su salida es un selector armónico. Aunque
aparecen los números \(12\) y \(-1\), no opera sobre series de Lambert ni
produce \(\epsR\).


\subsection{Origen dodecafásico del peso \(12:-1\) y comprobación diofántica de su coeficiente de polo}
% Fragmento público generado desde propietarios íntegros.
% Residencia: 52.8.3.
% Fuente: ../incoming/radion_monografia_20260827/manuscrito/sections/09_sector_90120.tex:115-175
\subsubsection{Funcional hexada--octada de Barbero--Immirzi}

\paragraph{Series orientadas y salto de incidencia}

Definimos
\begin{equation}
S_{90}(q)=\frac{q^{90}}{1-q^{270}},
\qquad
S_{120}(q)=\frac{q^{120}}{1-q^{360}},
\label{eq:barbero-series}
\end{equation}
y
\[
\Delta S_m=S_m(q_-)-S_m(q_+).
\]
El funcional de inclusión marcada es
\begin{equation}
K_{6\to8}=12\Delta S_{90}-\Delta S_{120}.
\label{eq:barbero-núcleo}
\end{equation}

\paragraph{Comprobación diofántica del peso dodecafásico \(12:-1\)}

La cadena fundamental dodecafásica y su lector de retorno, construidos en
\eqref{eq:l9s102:barbero-cadena-dodecafasica}--%
\eqref{eq:l9s102:barbero-peso-generado}, producen doce términos de tipo
\(S_{90}\) y un término de frontera de tipo \(-S_{120}\). El par orientado
queda determinado antes de calcular el coeficiente de polo.

Consideremos \(aS_{90}-bS_{120}\). La comprobación del coeficiente de
\((1-q)^{-1}\) toma la forma
\begin{equation}
\frac{a}{270}-\frac{b}{360}=\frac1{24}.
\label{eq:pole-condition}
\end{equation}
Multiplicar por \(1080\) da
\[
4a-3b=45.
\]
La única costura de la vuelta fundamental tiene multiplicidad absoluta uno.
Su verificación diofántica corresponde al contratérmino entero mínimo
\(|b|=1\). Para \(b=1\), \(a=12\); para \(b=-1\), \(a=21/2\) no es entero.
Por tanto,
\begin{equation}
\boxed{a:b=12:1,\qquad aS_{90}-bS_{120}=12S_{90}-S_{120}.}
\label{eq:12minus1}
\end{equation}

\begin{proposition}[Unicidad de la comprobación diofántica]
El par generado \((12,1)\) es la única solución entera de
\eqref{eq:pole-condition} con canal principal positivo y contratérmino
mínimo \(|b|=1\).
\end{proposition}

La realización hexada--octada conserva los canales \(90/120\)
previamente construidos; la cadena dodecafásica aporta las multiplicidades
y la orientación de frontera; el lector de retorno asigna las series a los
generadores. Estas operaciones conjuntas producen \(12S_{90}-S_{120}\).
La ecuación \eqref{eq:pole-condition}, la positividad y la minimalidad
conservan su función de comprobación aritmética posterior.

La cadena Barbero es posterior:
\[
(A,C^*)\to(x,y)\to q_\pm
\to K_{6\to8}\to\gamma_{\mathrm{BI}},
\]
con evaluación
\[
\gamma_{\mathrm{BI}}=0.274007672694459\ldots.
\]
Barbero--Immirzi no genera retroactivamente \(C^*\), \(\Delta_4\) ni
\(\epsR\).


\subsection{Extracción de la componente finita posterior al sistema dodecafásico orientado \(R_{12}\): \(C^*\mapsto C^*/6\), resto \(j\ge2\) y reconstrucción de \(6\) secciones}
% Fragmento público generado desde propietarios íntegros.
% Residencia: 52.8.4.
% Fuente: ../incoming/radion_monografia_20260827/manuscrito/sections/09_sector_90120.tex:31-54
\paragraph{Extracción de la componente \(C^*/6\) del registro dodecafásico}

Los doce sectores se agrupan en seis parejas opuestas:
\begin{equation}
p_i=e_i+e_{i+6},\qquad
a_i=e_i-e_{i+6},\qquad i=1,\ldots,6.
\label{eq:sexadic-pairs}
\end{equation}
La reconstrucción tiene perfil \(1+5+6=12\). La apertura por pareja es
\begin{equation}
h=\frac{C^*}{6}.
\label{eq:h-C6}
\end{equation}
No es una normalización exterior ni la sexta parte de un blanco: es el índice
interno de la reconstrucción sexádica del registro \(\Rstate\).

Definimos
\begin{equation}
q_{h,-}=\exp\left[-\frac\pi{180}(A-h)\right],
\qquad
q_{h,+}=\exp\left[-\frac\pi{180}(A+h)\right].
\label{eq:q-h}
\end{equation}


\subsection{Descomposición cuantitativa de las publicaciones de acarreo y espectral: \(F_{\mathrm{spec}}=\varepsilon_R^\circ+\chi_R\)}
% Fragmento público generado desde propietarios íntegros.
% Residencia: 52.8.5.
% Fuente: ../incoming/radion_monografia_20260827/manuscrito/sections/09_sector_90120.tex:55-114
\subsubsection{Publicación espectral de la cola de órdenes \(j\ge2\)}

\paragraph{Definición del funcional espectral}

Para \(m\in\{90,120\}\), sea
\begin{equation}
L_m^{(2)}(q)=\sum_{j\ge2}\frac{q^{mj}}{j}
=-\log(1-q^m)-q^m.
\label{eq:L2}
\end{equation}
La lectura de una pareja es
\begin{equation}
T_h^{(2)}=\frac{180}{\pi}
\Bigl[
L_{90}^{(2)}(q_{h,-})-L_{90}^{(2)}(q_{h,+})
-L_{120}^{(2)}(q_{h,-})+L_{120}^{(2)}(q_{h,+})
\Bigr].
\label{eq:T-h}
\end{equation}
La reconstrucción de las seis parejas da
\begin{equation}
F_{\mathrm{spec}}=6T_h^{(2)}
=5.1499235153094574410\ldots\times10^{-8}\,{}^\circ.
\label{eq:Fspec}
\end{equation}

El acarreo del radión es
\[
F_{\mathrm{acarreo}}=\epsR^\circ
=5.1383016758575282072\ldots\times10^{-8}\,{}^\circ.
\]
No coinciden:
\begin{equation}
\chi_R:=F_{\mathrm{spec}}-F_{\mathrm{acarreo}}
=1.16218394519292338\ldots\times10^{-10}\,{}^\circ.
\label{eq:chiR}
\end{equation}

La igualdad
\[
F_{\mathrm{acarreo}}=F_{\mathrm{spec}}-\chi_R
\]
es un cambio de carta válido una vez construidas ambas salidas. No debe
presentarse como generación independiente de \(\epsR^\circ\) si \(\chi_R\) se
define por la propia resta. La cola está cerrada y es significativa; lo falso
es identificarla directamente con el acarreo.

\begin{falsifier}
Tras convertir la cola expresada en grados a la unidad angular interna,
\[
\frac{\pi}{180^\circ}F_{\mathrm{spec}}-\epsone
=2.0283936357433839\ldots\times10^{-12}\ne0.
\]
Por tanto, ni el log completo ni la cola \(j\ge2\), aun después de aplicar la
carta de unidades, son \(\epsone\). Una
segunda factorización espectral autónoma necesitaría generar \(\chi_R\) sin
usar la diferencia como definición; el cierre APP--TPK del radión no depende
de esa promoción.
\end{falsifier}



\section{Holonomía orientada y geometría física: realizaciones posteriores de Minkowski, Hodge, Lorentz y Cartan--Holst}
% Fragmento público generado desde propietarios íntegros.
% Residencia: 52.9.
% Fuente: ../incoming/radion_monografia_20260827/manuscrito/sections/10_minkowski_hodge_lorentz.tex:1-55
\noindent La holonomía orientada admite realizaciones causales y geométricas posteriores en Minkowski, Hodge, Lorentz y Cartan--Holst.\par\medskip

\paragraph{Construcción de la métrica lorentziana desde el bloque APP}

\paragraph{Matriz APP de parámetros \(7\) y \(2\)}

La matriz central
\begin{equation}
B_c=\begin{pmatrix}7&2\\2&7\end{pmatrix}
=9P_++5P_-
\label{eq:Bc}
\end{equation}
se normaliza por su determinante:
\begin{equation}
M_c=\frac{B_c}{\sqrt{45}}
=\exp\left(\frac12\log\frac95\,R\right)
\in SO^+(1,1).
\label{eq:Mc}
\end{equation}
La rapidez interna es
\[
\eta_c=\frac12\log\frac95.
\]
La publicación proyectiva asociada satisface
\begin{equation}
P_B^n=P_++\left(\frac59\right)^nP_-.
\label{eq:PB-contract}
\end{equation}
La misma razón \(5/9\) que aparece como contracción espectral se reconoce aquí
como separación de un modo estable y otro contractivo. La identidad es
operatoria; no convierte todas las apariciones de \(5/9\) en el mismo mapa.

\paragraph{Incidencia \(6\to8\) y cociente de firma causal}

Sea \(M\) la incidencia tipada entre seis caras y ocho vértices. La relación
espectral recuperada es
\begin{equation}
MM^{\mathsf T}=12P_u+4P_-,
\label{eq:MMt}
\end{equation}
donde \(P_u\) proyecta sobre el modo uniforme y \(P_-\) sobre el subespacio
orientado. El cociente superviviente de dimensión cuatro se descompone como
\begin{equation}
Q_4\simeq\R u\oplus E_-.
\label{eq:radion-Q4}
\end{equation}
Declarando negativo el modo uniforme y positivos los tres modos orientados,
se obtiene
\begin{equation}
B_{\mathrm{caus}}=\operatorname{diag}(-1,1,1,1).
\label{eq:minkowski}
\end{equation}
La pantalla de Minkowski no selecciona el estado HMT. Es la realización
causal posterior del cociente que la incidencia ya construyó.


\subsection{Dualidad de Hodge sobre \(2\)-formas lorentzianas}
% Fragmento público generado desde propietarios íntegros.
% Residencia: 52.9.1.
% Fuente: ../incoming/radion_monografia_20260827/manuscrito/sections/10_minkowski_hodge_lorentz.tex:56-100
\subsubsection{Dualidad de Hodge y operador de cuarto de giro}

Fijadas la firma \((-+++ )\) y una orientación, la estrella de Hodge sobre
dos-formas satisface
\begin{equation}
\star:\Lambda^2Q_4^*\to\Lambda^2Q_4^*,
\qquad \star^2=-I.
\label{eq:hodge-star}
\end{equation}
De este modo, el cuarto de giro que en el plano de fase se expresaba mediante
\(J^2=-I\) reaparece en el espacio de campos como estructura compleja sobre
\(\Lambda^2\). Los dominios son diferentes, pero la genealogía de orientación
se conserva mediante el funtor.

Sea \(A_{\mathrm{em}}\) un potencial y
\begin{equation}
F=\diff A_{\mathrm{em}}.
\label{eq:F-dA}
\end{equation}
La descomposición de \(F\) respecto a un observador temporal produce campos
eléctrico y magnético; \(\star F\) intercambia sus papeles con el signo fijado
por la orientación y la firma. La polaridad constitutiva de
\eqref{eq:EM-conjugation} adquiere así una realización diferencial.

\paragraph{Cadena de transducción electromagnética y causal}

El transporte que debe mantenerse visible es
\begin{equation}
\Gamma_9
\longrightarrow\text{hélice de memoria}
\longrightarrow(Q_4,B_{\mathrm{caus}},\star)
\longrightarrow A_{\mathrm{em}}
\longrightarrow F=\diff A_{\mathrm{em}}
\longrightarrow
m\nabla_u u=q(\iota_uF)^\sharp.
\label{eq:lorentz-chain}
\end{equation}
Cada flecha añade estructura: la holonomía conserva retorno; el cociente fija
causalidad; Hodge fija dualidad; el potencial fija el campo; el acoplamiento
con carga produce dinámica.

La ecuación final es la fuerza de Lorentz en forma geométrica. La hélice TPK
es su antecedente de memoria, no una órbita física ya movida por un campo. La
dinámica aparece sólo al introducir \(F\), \(q\) y la conexión \(\nabla\).


\subsection{Separación entre memoria torsional interna y realización geométrica de Cartan--Holst}
% Fragmento público generado desde propietarios íntegros.
% Residencia: 52.9.2.
% Fuente: ../incoming/radion_monografia_20260827/manuscrito/sections/10_minkowski_hodge_lorentz.tex:101-186
\subsubsection{Realización lorentziana de las trayectorias helicoidales}

En un campo magnético uniforme, una partícula de carga \(q\), masa \(m\) y
factor relativista \(\gamma\) posee
\begin{equation}
\omega_c=\frac{|q|B}{\gamma m},
\qquad
r_L=\frac{p_\perp}{|q|B},
\qquad
\Delta z=\frac{2\pi p_\parallel}{|q|B}.
\label{eq:lorentz-helix}
\end{equation}
La fase ciclotrónica cierra mientras el movimiento paralelo avanza: la órbita
es helicoidal. El paralelismo con \eqref{eq:H9-helix} es estructural y puede
formalizarse por un transductor que preserve fase, orientación y paso. No se
identifica el cociente \(q_9\) con la coordenada espacial \(z\); se reconoce
la misma arquitectura de retorno más avance.

La fuerza de Lorentz no se deriva de la imagen de una espiral por semejanza.
Se deriva al completar \eqref{eq:lorentz-chain}. La imagen ayuda a comprender
por qué fase circular y memoria longitudinal son compatibles; la forma de
campo determina la dinámica.

\subsubsection{Confluencia entre acción, área, información y gravedad}

\paragraph{Realización Cartan--Holst de la holonomía}

La promoción gravitatoria se expresa mediante un morfismo de holonomías
\begin{equation}
\operatorname{Hol}_{\mathcal A}
\bigl(\mathfrak R_{\mathrm{grav}}(\gamma)\bigr)
=\rho_C\bigl(\operatorname{Hol}_{\TPK}(\gamma)\bigr).
\label{eq:holonomy-gravity}
\end{equation}
El lado derecho contiene la memoria de ruta; \(\rho_C\) la realiza en una
conexión gravitatoria. La torsión de Cartan
\[
T^I=D_\omega e^I
\]
aparece en el codominio. La torsión global \(\Delta_4\) es una coordenada
upstream que puede alimentar el mapa, pero no se renombra como \(T^I\) sin
\eqref{eq:holonomy-gravity}.

\paragraph{Transducción del área de acción en área informacional}

El radión fija un cuanto de acción por fase. El lector hexada--octada de
Barbero usa posteriormente los canales \(q_\pm\), la incidencia \(90/120\) y
el funcional \eqref{eq:barbero-núcleo}. La dirección causal es
\[
(A,C^*,\hret)\to q_\pm\to K_{6\to8}
\to\gamma_{\mathrm{BI}}\to\text{área/información}.
\]
El alfabeto trítico aporta la unidad de entropía \(\log3\) y un corte de
\(81\) canales aporta \(81\log3\). Esta entropía no se confunde con
\(\log2\), con \(\log\varphi\), con una entropía microcanónica ni con una
entropía de von Neumann sin especificar el estado.

La síntesis física es afirmativa pero tipada: el radión aporta la unidad de
acción orientada que una holonomía puede transportar; Barbero publica la
normalización de área; la pantalla causal y Hodge permiten realizar campos;
Cartan--Holst recibe la holonomía en geometría gravitatoria.

\begin{statusbox}
La ecuación escalar del radión no es por sí sola una ecuación de Einstein. Su
aportación a la gravedad es ser una sección exacta de la arquitectura de
acción, orientación y holonomía que la realización gravitatoria consume.
\end{statusbox}

\subsubsection{Transformación de Wick como cambio de lector analítico}

En la formulación convencional, la rotación de Wick relaciona una variable
temporal lorentziana con una coordenada euclídea mediante una continuación
compleja. En HMT, la raíz \(J^2=-I\), el involutivo \(R^2=I\) y el umbral
\(N^2=0\) ya existen antes. El cambio de lector
\[
J\longleftrightarrow R
\]
transporta una pareja común/orientada entre publicación elíptica y
publicación hiperbólica. La continuación compleja reconoce después esa
operación en coordenadas analíticas.

Esto explica por qué círculo, elipse de acción e interior hiperbólico pueden
pertenecer al mismo estado sin ser la misma métrica. El círculo usa \(J\) para
la fase; la forma de la elipse usa un compresión simpléctica generado por \(R\); Klein mide
el interior mediante razón doble; la pantalla de Minkowski realiza causalidad
en otro codominio.


\section{Elipse, retículo y fase de camino: carácter modular, órbitas y propagación}
\subsection{Red rectangular de períodos, involución \(\tau\mapsto-1/\tau\) e invariancia de \(j\)}
% Fragmento público generado desde propietarios íntegros.
% Residencia: 52.10.1.
% Fuente: ../incoming/radion_monografia_20260827/manuscrito/sections/11_modularidad_orbitas.tex:1-84
\noindent Los caracteres modulares, orbitales y de camino conservan invariantes distintos de una misma forma orientada.\par\medskip

\subsubsection{Toro complejo asociado a la elipse de acción}

\paragraph{Retículo rectangular de períodos}

La elipse por sí sola no es un toro. Declaramos el retículo
\begin{equation}
\Lambda_\eta=a_S\Z+i,b_S\Z
\label{eq:lattice}
\end{equation}
y el cociente \(\C/\Lambda_\eta\). Su módulo es
\begin{equation}
\tau_\eta=i\frac{b_S}{a_S}=ie^{-\eta}
=i\,0.741048144159375160795483663369300\ldots.
\label{eq:tau-eta}
\end{equation}
La inversión bidireccional \(\eta\mapsto-\eta\) intercambia los ejes:
\begin{equation}
\tau_{-\eta}=ie^\eta=-\frac1{\tau_\eta}.
\label{eq:modular-S}
\end{equation}
Es exactamente la transformación modular \(S:\tau\mapsto-1/\tau\).

\paragraph{Invariancia de la función modular \(j\)}

Por invariancia modular,
\begin{equation}
\boxed{j(\tau_{-\eta})=j(\tau_\eta).}
\label{eq:j-invariance}
\end{equation}
En la normalización \(j(i)=1728\),
\[
j(\tau_\eta)
=5597.43843932408729830097089254768\ldots.
\]
La función moonshine normalizada \(J=j-744\) da
\[
J(\tau_\eta)
=4853.43843932408729830097089254768\ldots.
\]
Para el bloque APP basal,
\[
\tau_0=i\frac{\sqrt5}{3},
\qquad
j(\tau_0)=5369.48832024674306021916882626352\ldots.
\]

\begin{proposition}[Lo que conserva la clase modular]
La involución de hoja invierte \(\eta\), pero \(j\) conserva la clase del
toro. Los focos y el etiquetado de ejes retienen la orientación que \(j\)
olvida.
\end{proposition}

Este resultado conecta de manera exacta la doble vía con la modularidad una
vez declarado el cociente. No autoriza a hacer actuar directamente al
Monstruo sobre la carga, los dígitos o los focos. El corredor excepcional
\[
\mathrm{APP}\to\mathrm{Paley}\to\mathrm{Witt}\to\mathrm{Golay}
\to A_2^{12}\to\mathrm{Leech}\to V^\natural
\to\mathrm{Monster}\to\mathrm{Moonshine}
\]
es otra proyección del mismo estado enriquecido. Para identificar una acción
concreta sobre el radión se necesita un entrelazador que preserve sus
coordenadas pertinentes.

\paragraph{Cadena de invariantes del carácter de forma}

Reuniendo las realizaciones declaradas,
\begin{equation}
\boxed{
r_\eta=e^{-\eta}
=\sqrt{\frac{A-C^*}{A+C^*}}
=\frac{b_1}{a_1}
=\frac{b_S}{a_S}
=\frac{\Delta P}{\Delta Q}
=\frac{Z_\eta}{Z_*}
=\Im\tau_\eta.}
\label{eq:shape-master}
\end{equation}
Una sola razón recorre geometría espectral, acción, covarianza, impedancia y
módulo. Cada igualdad se apoya en un mapa explícito; ninguna depende de una
semejanza decimal.


\subsection{Compatibilidad entre el carácter modular y la propagación sobre rutas orientadas}
% Fragmento público generado desde propietarios íntegros.
% Residencia: 52.10.2.
% Fuente: ../incoming/radion_monografia_20260827/manuscrito/sections/11_modularidad_orbitas.tex:85-187
\subsubsection{Conservación areal en las órbitas de Kepler y Sommerfeld}

\paragraph{Equivalencia areal de las realizaciones orbital y oscilatoria}

La elipse de Kepler vive en espacio de configuración y responde a un potencial
central. La elipse del radión vive en espacio de fases y responde a una forma
de acción. Son objetos diferentes. Comparten una ley areal porque ambas
realizaciones poseen una dos-forma conservada.

Para una órbita kepleriana,
\begin{equation}
r(\phi)=\frac{p}{1+e\cos(\phi-\phi_0)},
\qquad
\frac{\diff A}{\diff t}=\frac{\ell}{2m}.
\label{eq:kepler}
\end{equation}
La tercera ley toma la forma
\[
T^2=\frac{4\pi^2\mu}{K}a^3.
\]
Para el oscilador en fase,
\begin{equation}
J=\oint p\,\diff q=\frac{2\pi E}{\omega}.
\label{eq:osc-action}
\end{equation}
Si la celda elemental es \(J=\hfull\), entonces
\begin{equation}
E=\hret\omega,\qquad \frac{J}{2\pi}=\hret.
\label{eq:E-hbarw}
\end{equation}
Ésta es la convergencia limpia entre acción de Planck, frecuencia angular y
radión.

Sommerfeld cuantiza integrales de acción sobre ciclos. La lectura HMT añade
la memoria de qué ciclo, hoja y retorno sobreviven. En una banda de Möbius,
un bucle visible puede portar \(\hfull/2\) y dos bucles restaurar \(\hfull\).
Ese factor topológico no se suma al índice de Maslov de una cuantización EBK;
pertenecen a correcciones con propietarios distintos.

\paragraph{Holonomía orbital y propagaciones avanzada y retardada}

Una órbita que acumula holonomía satisface una condición del tipo
\begin{equation}
\epsilon_\gamma\exp\left(\frac{iJ_\gamma}{\hret}\right)=1.
\label{eq:holonomy-action}
\end{equation}
El signo \(\epsilon_\gamma\) conserva la hoja; \(J_\gamma\) conserva la
acción. El adelanto o retardo apsidal se interpreta como desacoplamiento entre
el cierre angular visible y el retorno completo. El radión suministra la
unidad orientada con la que se mide esa diferencia.

\subsubsection{Composición de amplitudes de Schrödinger y suma de caminos}

\paragraph{Cuatro realizaciones de una misma fase}

Una vez fijado \(\hret\), distintas ecuaciones convencionales reconocen
aspectos del mismo transporte:
\begin{enumerate}
\item Schrödinger publica la fase mediante
\(i\hret\partial_t\psi=H\psi\).
\item Pauli añade la doble hoja de espín y su acoplamiento magnético.
\item Dirac linealiza la relación energía--momento y organiza partícula y
antipartícula en una representación causal.
\item La suma de caminos asigna a cada historia la fase
\(\exp(iS[\gamma]/\hret)\).
\end{enumerate}
Estas teorías no generan retrospectivamente \(\hret\). Reciben la unidad de
acción y realizan diferentes contenidos del estado enriquecido.

\paragraph{Funcional de fase sobre el espacio de caminos}

En la formulación de Feynman,
\begin{equation}
\mathcal K(b,a)=\int_{\gamma:a\to b}
\exp\left(\frac{i}{\hret}S[\gamma]\right)\mathcal D\gamma.
\label{eq:path-integral}
\end{equation}
Un incremento de acción \(\hret\) cambia la fase en un radián. Un incremento
\(\hfull\) la cambia en una vuelta. El teorema
\(\mathcal A(\theta)=\hret\theta\) convierte esta relación en geometría de
área: la fase de un camino puede leerse como acción barrida en unidades de
radión.

La inversión \(C_M=-\operatorname{rev}\) sobre palabras dodecafásicas y la
identidad \eqref{eq:time-reversal} organizan rutas conjugadas. En una
realización de absorbedor se separan
\[
G_{\mathrm{sym}}=\tfrac12(G_{\mathrm{ret}}+G_{\mathrm{adv}}),
\qquad
G_{\mathrm{odd}}=G_{\mathrm{ret}}-G_{\mathrm{adv}}.
\]
El primer sector es par bajo intercambio retardado--avanzado; el segundo es
impar. Esto reproduce la separación entre energía/forma par y acción
orientada impar del radión cuando el transductor preserva frontera y forma
simpléctica.

\begin{statusbox}
La teoría clásica del absorbedor y la interpretación Feynman--Stückelberg de
antipartículas son realizaciones históricas diferentes. La primera usa una
interacción temporalmente simétrica; la segunda reorganiza propagación en la
teoría relativista. HMT aporta un antecedente común de reversión, pero no las
funde sin un mapa de propagadores.
\end{statusbox}


\section{Radical nonádico y frontera decimal: conservación del acarreo del radión}
\subsection{El acarreo del radión en la completación decimal}
% Fragmento público generado desde propietarios íntegros.
% Residencia: 52.11.1.
% Fuente: ../incoming/radion_monografia_20260827/manuscrito/sections/12_aritmetica_borde.tex:1-159
\noindent La frontera nonádica conserva el defecto suma--producto, la vacancia y el acarreo de completación.\par\medskip

\subsubsection{Radical \(3\! -\! 6\! -\! 9\) del borde}

\paragraph{Generación y nilpotencia del radical nonádico}

En el anillo \(\Z/9\Z\), el ideal
\begin{equation}
\mathfrak m=(3)=\{0,3,6\}
\label{eq:radical-369}
\end{equation}
satisface
\begin{equation}
\mathfrak m^2=(9)=0.
\label{eq:m-square-zero}
\end{equation}
Es el sector nilpotente del borde nonádico. La órbita visible de sus tres
marcas se publica como
\[
369\longrightarrow693\longrightarrow936\longrightarrow369.
\]
El retorno no implica que todos los estados sean iguales: la permutación
visible cierra mientras los cociclos de ruta pueden avanzar.

La doble convención
\[
9\equiv0\pmod9,
\qquad
\operatorname{dr}_9(9)=9
\]
expresa dos lectores: cero residual algebraico y cierre visible de raíz
digital. Por eso un nueve puede absorber torsión de borde sin convertirse en
un cero específico de la función zeta.

\paragraph{Defecto suma--producto \(459-405=54\)}

Las hojas APP producen totales complementarios
\[
S_+=405,\qquad S_\times=459,\qquad S_\times-S_+=54.
\]
El \(54\) es un defecto exacto de las tablas, y reaparece como exponente de
contracción por la identidad \(9\cdot6=54\). Sin embargo, cada aparición
requiere su mapa. La distancia focal \(d_1=0.544545\ldots\) no es igual a
\(0.54\), y el análisis de invariantes no ha producido una flecha que la
identifique con el defecto \(54\).

La potencia
\[
9^9=387,420,489
\]
genera una escala donde
\begin{equation}
\operatorname{ord}_{10}\bigl(9^{9^9}\bigr)=369,693,099,
\qquad
\#\,\text{dígitos}=369,693,100.
\label{eq:369693}
\end{equation}
El bloque \(369|693\) aparece en el par orden--longitud; no se afirma que sea
el prefijo decimal de la potencia.

\subsubsection{Teoría aritmética de las vacancias decimales}

\paragraph{Pendiente asintótica de vacancia}

La separación entre la década y la base nonádica es
\begin{equation}
\delta_{\mathrm{vac}}=\log_{10}\frac{10}{9}.
\label{eq:delta-vac}
\end{equation}
El contador acumulado
\begin{equation}
V_9(n)=\left\lceil n\delta_{\mathrm{vac}}\right\rceil
\label{eq:vacancy-count}
\end{equation}
registra cuántas posiciones de corrección necesita la publicación decimal de
una evolución nonádica. Sus saltos forman una palabra binaria
\[
z_n=V_9(n)-V_9(n-1)\in\{0,1\}.
\]
Los huecos se distribuyen con longitudes \(21/22\), los retornos con
longitudes \(6/7\), y el determinante combinado es \(15\). Estos números
proceden del lector de vacancia, no del núcleo Barbero aunque el \(15\)
coincida con \(\binom62\).

\paragraph{Polarización y corriente aritmética de la red nonádica}

Definimos
\begin{equation}
P_N=\sum_{n=1}^N(z_n-\delta_{\mathrm{vac}}),
\qquad
J_N=P_N-P_{N-1}=z_N-\delta_{\mathrm{vac}}.
\label{eq:arith-current}
\end{equation}
La secuencia equilibrada satisface \(|P_N|<1\). La \enquote{corriente}
\(J_N\) es aquí una corriente aritmética: mide creación y absorción de
vacancias respecto a la pendiente media. Sólo un transductor dimensional
posterior puede publicarla como densidad de corriente electromagnética.

La intuición física es potente porque la misma forma algebraica separa un
promedio par y una fluctuación orientada. El rigor exige conservar el tipo:
\[
J_N\in\R\quad\text{adimensional},
\qquad
j^\mu\in\text{recta física de corriente}.
\]

\subsubsection{Completación \(729\to1000\) y conservación del acarreo}

\paragraph{Descomposición \(1000=729+243+27+1\)}

La completación
\[
1000=729+243+27+1
\]
conserva la unidad mediante cuatro escalas ternarias. La corona total es
\[
1000-729=271,
\]
mientras el último entero antes del cierre satisface
\[
999=729+270.
\]
La diferencia entre \(270\) y \(271\) distingue frontera ocupada de unidad
completante. El mismo cuidado evita identificar una vacancia con un cero o un
acarreo con un residuo real.

\paragraph{Residencia del acarreo del radión en la corona decimal}

Los operandos del radión comienzan, en bloques de tres cifras, por
\[
\Phi_T:\quad000|027|455|906|837|565|\cdots,
\]
\[
D_E:\quad000|027|455|010|034|743|\cdots.
\]
La resta APP da
\[
\epsone:\quad000|000|000|896|802|822|\cdots.
\]
El pequeño tamaño del resultado no se obtiene truncando los sumandos bajos.
Se obtiene porque dos publicaciones distintas comparten un largo prefijo y el
operador \(\operatorname{SubAcarreo}_{1000}\) conserva los préstamos necesarios
para restarlas exactamente.

La palabra \emph{memoria} tiene aquí dos niveles:
\begin{enumerate}
\item \(q_{\mathrm{mem}}\) registra la profundidad de retornos nonádicos;
\item \(c_i\) registra los préstamos locales de la resta decimal.
\end{enumerate}
El primero organiza la hélice; el segundo publica el decimal. \(\epsR\) usa
ambos porque sus operandos proceden de cilindros a profundidad y su resta se
realiza con acarreo, pero no es igual a ninguno de los contadores.

\begin{thesisbox}
El borde no es un lugar donde la teoría pierde información. Es el lugar donde
la información cambia de representación: fase que vuelve, memoria que
asciende, cilindro que se contrae, vacancia que se registra y acarreo que
transporta la diferencia.
\end{thesisbox}

\subsection{Compatibilidad de las secciones directa y torsional: caracterización estructural de \(\pi\)}
% Fragmento público generado desde propietarios íntegros.
% Residencia: 52.11.2.
% Fuente: ../incoming/radion_monografia_20260827/manuscrito/sections/03_cierre_exacto.tex:286-309
\subsubsection{Caracterización estructural de \(\pi\) por compatibilidad de secciones}

Reordenar \eqref{eq:cierre-unitario} da
\begin{equation}
\boxed{
2\pih=
\frac{1+\Phi_T-\epsone}
{\frac5{36}+\frac{25}{9}\alphah}.}
\label{eq:pi-radion}
\end{equation}
Esta igualdad puede llamarse una nueva definición de \(\pi\) dentro de la
carta del radión si se entiende \emph{definición estructural} como una
caracterización exacta por compatibilidad de secciones. No es un segundo
generador independiente: \(\Delta_4\) y el operador de acarreo ya reciben la
sección coinductiva de \(\pi\). La novedad está en que la vuelta queda
caracterizada por el empalme entre publicación directa, torsión y memoria.

\begin{narrativebox}
La ecuación \eqref{eq:pi-radion} no dice que primero desconozcamos \(\pi\) y
lo adivinemos cerrando un ángulo. Dice que el \(\pi\) generado por la conexión
nonádica es exactamente el que hace compatibles dos lecturas internas del
mismo estado. Es una identidad de reconocimiento intrínseco, no una
calibración metrológica.
\end{narrativebox}


\section{Conservación evolutiva de la unidad holográfica a través de las dimensiones}
% Conversión TeX fiel del cuerpo científico del delta narrativo.
% Fuente congelada: ../../../../../HMT2/CONTINUIDAD_EDITORIAL/RECONSTRUCCION_ARMONICA_MACROTRATADO_20260822/REVISION_CONJUNTA_PARTES_I_II/auditar_pdf_rigor_academico/docs/DELTA_NARRATIVO_UNIDAD_RADION_AUTOCONSERVACION_AUTOINERCIA_GRAVEDAD_PENTADICA_20260828.md:12-366.
\subsection{Separación tipada entre el régimen de curvatura \(\tau\in\{+1,0,-1\}\), la respuesta radial \(p\) y la holonomía global de espín}

La biografía del radión puede representarse mediante una curva elevada

\[
\gamma(\theta)
=
\bigl(r(\theta)\cos\theta,\,
r(\theta)\sin\theta,\,
m(\theta)\bigr),
\]

donde las dos primeras coordenadas pertenecen al plano euclídeo \(\mathbb R^2\) y \(m(\theta)\) registra la memoria o profundidad del levantamiento. Esta ecuación permite leer conjuntamente círculo, hélice y espiral.

Cuando evoluciona el ángulo y permanecen constantes el radio y la memoria, la publicación plana es un círculo. Cuando evolucionan ángulo y memoria mientras el radio conserva su valor, la trayectoria completa es una hélice de radio constante. Cuando evolucionan ángulo y radio sobre una memoria fija, aparece una espiral en \(\mathbb R^2\). Cuando evolucionan las tres coordenadas, la trayectoria es una hélice radial: una biografía en la que fase, amplitud y memoria se transforman conjuntamente.

El carácter radial \(p\) clasifica la ley de respuesta del radio. Una familia útil es

\[
r_p(\theta)
=
r_0\bigl[1+p\lambda(\theta-\theta_0)\bigr]^{1/p},
\qquad p\neq 0,
\]

con límite logarítmico

\[
r_0\exp\bigl(\lambda(\theta-\theta_0)\bigr)
\qquad\text{cuando }p\to0.
\]

En esta parametrización, \(p=2\) publica el régimen de Fermat, \(p=1\) el régimen arquimediano, \(p=0\) el logarítmico, \(p=-1\) el hiperbólico y \(p=-2\) el de lituus, con las elecciones de dominio y orientación correspondientes. Los valores fraccionarios describen regímenes intermedios. Así, \(p=\tfrac12\), \(p=\tfrac32\) o \(p=-\tfrac12\) expresan respuestas radiales que interpolan entre las familias enteras y amplían la clasificación geométrica disponible.

El carácter también puede cambiar al atravesar un umbral. La escritura diferencial

\[
\frac{dx}{d\theta}
=
\beta(m)\,x^{1-p(m)},
\qquad x=\frac r{r_0},
\]

describe una biografía por regímenes: \(p(m)\) permanece estable dentro de una fase y cambia cuando la memoria alcanza una frontera de publicación. La trayectoria resultante enlaza sectores circulares, helicoidales y espirales mediante una ley común. La continuidad del estado y el registro del umbral convierten esos sectores en fases de una misma evolución.

Tres invariantes ocupan aquí niveles distintos y coordinados. El signo de curvatura

\[
\tau\in\{+1,0,-1\}
\]

clasifica los regímenes elíptico, parabólico e hiperbólico. El carácter \(p\) clasifica la respuesta radial. El espín registra la holonomía de orientación acumulada por el levantamiento. Curvatura, respuesta radial y espín forman un sistema de lectura: la curvatura determina el tipo local de espacio, \(p\) determina cómo la trayectoria ocupa ese espacio y el espín conserva la orientación adquirida al recorrerlo. Las curvas cerradas y abiertas aparecen entonces como biografías geométricas distintas de la misma unidad orientada.

\subsection{Partícula como sección persistente de una órbita enriquecida: fase, ruta, fibra, multisección y memoria}

La Mecánica Dimensional puede formular la partícula como una sección persistente de esa extensión geométrica y dinámica. Una notación conceptual compacta es

\[
\mathfrak P
=
\operatorname{Sec}_{\mathrm{pers}}
\bigl[
\tau,p,\operatorname{Hol};
\text{hoja},\text{fibra},\text{ruta},\text{memoria}
\bigr].
\]

La partícula conserva una ruta reconocible a través de cambios de fase, curvatura y dimensión. Su espín publica la holonomía de orientación; su carga publica la polaridad de fase; su color publica una orientación interna de la fibra calibre; su sabor distingue familias de rutas o multisecciones; su masa expresa el carácter de transducción dimensional mediante el cual la acción areal adquiere lectura volumétrica. Cada propiedad es un lector del estado persistente y todas remiten a una genealogía común.

El electrón ocupa el umbral basal de esta construcción porque ofrece el primer lector estable en el que fase, carga, espín y masa quedan simultáneamente publicadas. La masa electrónica rectora

\[
m_e^\beta=0.510998950690000808608\ldots\ \mathrm{MeV}
\]

fija una referencia de transducción; los estados basales anteriores conservan su función como etapas genealógicas. A partir del electrón, los cocientes \(K_D/K_\Omega\), las torres \(90/120\) y el dominio \(81/56/13/324\) organizan la publicación de familias masivas. La ley de masas lee, por tanto, cómo cada sección persistente atraviesa el espacio enriquecido y convierte densidad areal de acción en presencia volumétrica.

Esta interpretación une las espirales con las partículas de manera precisa. El índice \(p\) aporta la respuesta radial; \(\tau\) aporta el régimen de curvatura; la holonomía aporta el espín; la ruta y la fibra aportan sabor y color; la transducción área–volumen aporta el carácter masivo. El atlas de partículas se convierte así en una clasificación de modos estables de atravesamiento dimensional.

\subsection{Transducción de la fase orientada en polarización eléctrica, respuesta magnética y emisión radiativa de frontera}

La electricidad se presenta como transporte dirigido de fase y acción a través de una estructura de vacancias. Una vacancia es una transición admisible del soporte: abre un canal, fija un umbral y permite que la actualización avance. La componente eléctrica publica el transporte tangencial a lo largo de la ruta. La componente magnética publica la curvatura transversal inducida por ese transporte. Su perpendicularidad expresa la relación geométrica entre avance y giro.

Cuando dos componentes transversales sinusoidales mantienen el desfase adecuado, el vector del campo gira. La propagación longitudinal eleva ese giro a una hélice. La variación de amplitud transforma su proyección frontal en una espiral. La lectura geométrica queda entonces completamente coordinada: el ángulo porta la fase; la velocidad angular porta la frecuencia; el radio porta amplitud y densidad de acción; la orientación porta polarización y sentido; la coordenada vertical porta propagación y memoria; \(p\) porta el régimen radial; el umbral porta emisión, absorción o cambio de estado.

El electromagnetismo aparece como una dinámica de fase–acción–memoria con publicaciones complementarias. La electricidad registra el avance polar; el magnetismo registra la respuesta transversal; la onda registra la periodicidad; la hélice registra el avance con memoria; la espiral registra la evolución radial; la radiación registra la publicación del estado cuando la acumulación alcanza un umbral.

Esta arquitectura también ordena el papel de las vacancias y los ciclos. La vacancia proporciona la posibilidad local de transición; el ciclo conserva la fase; la memoria distingue las vueltas; el umbral decide la residencia siguiente de la acción. La propagación puede permanecer confinada, cambiar de régimen radial o publicarse como radiación controlada. Las fuerzas de campo adquieren así una descripción genealógica continua desde el soporte discreto hasta la onda observada.

\subsection{Reversibilidad de la dinámica completa y coste termodinámico de las proyecciones reductoras de información}
El Landauer nulo constituye un caso límite de esta ley. Para una actualización biyectiva del estado completo,

\[
H(X\mid U(X))=0,
\qquad
\inf Q_L=0.
\]

La información permanece disponible en la genealogía. Cuando un lector realiza un cociente \(q\), la entropía se descompone como

\[
H(X)=H(q(X))+H(X\mid q(X)).
\]

El término condicional mide la información retenida fuera de la proyección. Un reinicio trítico materializado satisface el umbral

\[
Q_{\mathrm{reset}}\ge k_B\Theta\log 3.
\]

La constante de Boltzmann traduce la multiplicidad que el lector térmico agrupa; el umbral de Nyquist separa el muestreo fiel del aliasing; el Landauer nulo caracteriza el transporte reversible del estado íntegro. Información, temperatura y acción quedan articuladas por el tipo de lectura aplicado.

\subsection{Fuerzas autoconservativas como transporte de acción entre los sectores observable, radial, de memoria, de vacancia y de radiación}

La conservación evolutiva se expresa mejor mediante intercambio que mediante inmovilidad escalar. Un estado completo distribuye acción y momento entre movimiento visible, memoria, deformación radial, vacancias, radiación y rama conjugada de retorno. El balance enriquecido puede escribirse como

\[
\Delta J_{\mathrm{visible}}
+\Delta J_{\mathrm{memoria}}
+\Delta J_{\mathrm{radial}}
+\Delta J_{\mathrm{vacancias}}
+\Delta J_{\text{radiación}}
+\Delta J_{\mathrm{retorno}}
=0.
\]

Una fuerza es \textbf{autoconservativa} cuando su propia ley de actualización contiene los canales que transportan, almacenan, publican y reabsorben la acción. La conservación pertenece al balance genealógico completo de la interacción. Cada lector observa una parte del intercambio; la genealogía recompone su totalidad.

El sistema Tierra–Luna ofrece una imagen física directa. La rotación terrestre, la órbita lunar, las mareas, la deformación y la radiación intercambian momento angular. El consumo neto del sistema completo se cierra en el balance genealógico de esas transferencias. La magnitud denominada energía funciona como lectura derivada de la acción y del momento publicados en cada sector. La conservación fundamental reside en la compatibilidad entre todos los lectores del intercambio.

Esta formulación disuelve la antigua frontera entre fuerzas conservativas y fuerzas llamadas no conservativas mediante una categoría más profunda. Rozamiento, disipación aparente y radiación ocupan canales explícitos del estado ampliado. Una pérdida en la proyección mecánica visible corresponde a una transferencia hacia memoria, deformación, calor, vacancias o radiación. La autoconservación sigue la residencia de la acción a través de esos cambios.

\subsection{Sistemas autoinerciales: publicación de la aceleración mediante orientación, curvatura, hoja y memoria}

La expresión adecuada es \textbf{autoinercia relacional}. Designa la compatibilidad entre la conexión propia del estado total y la respuesta que aparece al proyectarlo dentro de la distribución global. Los aspectos endógeno y exógeno son componentes analíticas de una sola relación dinámica.

Sea \(\Gamma\) una trayectoria horizontal o geodésica del estado completo. Su ley propia satisface

\[
\nabla_{\dot\Gamma}\dot\Gamma=0.
\]

Al proyectarla sobre un lector \(S\), la aceleración observada queda determinada por la curvatura de la proyección:

\[
\nabla^S_{\dot\Gamma_S}\dot\Gamma_S
=
\mathcal K_{\pi_S}(\dot\Gamma,\dot\Gamma).
\]

La aceleración publicada expresa la relación entre movimiento total y geometría del lector. El principio de Mach completa este mecanismo mediante la traza global de frontera

\[
b:\mathcal X_\infty\longrightarrow B_\infty
\]

y la factorización de la respuesta inercial local

\[
I_v=\bar I_v\circ b.
\]

La distribución global entra en la respuesta local a través de una aplicación tipada. La autoinercia relacional reúne, por tanto, la conexión endógena del movimiento, la condición global transmitida por \(b\) y la proyección concreta del observador. Una escritura sintética para esa coordinación es

\[
\nabla^{\mathrm{auto}}
=
\mathcal M\bigl(
\nabla^{\mathrm{end}},
b[\mathcal X_\infty],
\pi_S
\bigr),
\]

donde \(\mathcal M\) representa el acoplamiento machiano y conserva el tipado de cada componente.

El principio de Ambivalencia añade la doble orientación de la lectura. Toda publicación directa conserva una rama conjugada de retorno; emisión y absorción forman extremos relacionados de la misma biografía; la orientación temporal se lee desde ambos sentidos del transporte. La teoría del absorbedor encuentra aquí una formulación natural: fuente y receptor cierran conjuntamente el balance genealógico de fase y acción. El presente torsional es la interfaz donde ambas orientaciones se compatibilizan.

TRIT actúa en este nivel como cristal de tiempo: organiza la periodicidad local de la actualización y conserva, mediante memoria, la diferencia entre retornos sucesivos. La rama directa y la rama conjugada relacionan emisión con absorción, y el acoplamiento observador–observable adquiere una residencia dinámica dentro del mismo estado. El tiempo publicado es la lectura de esa orientación; el tiempo genealógico es la secuencia completa de fase, retorno y memoria.

Mach y Ambivalencia cumplen funciones complementarias. Mach relaciona la respuesta local con la distribución global. Ambivalencia conserva las direcciones directa y conjugada de la evolución. La autoinercia relacional articula ambas y sustituye la escisión entre sistemas inerciales y no inerciales por una ley única de conexión, proyección y memoria.

\subsection{Dualidad areal--volumétrica, transición hexada--octada y función del parámetro de Barbero--Immirzi}

El paso de área a volumen constituye una transducción dimensional central. El área publica acción de frontera; el volumen publica profundidad interior. La derivación HMT del parámetro de Barbero–Immirzi organiza la compatibilidad entre ambas medidas mediante la incidencia hexada–octada y los canales \(90/120\):

\[
90=6\binom62,
\qquad
120=8\binom62.
\]

Los cardinales \(6\) y \(8\) admiten una realización cúbica mediante caras y vértices, mientras la hexada y la octada de Steiner constituyen objetos de incidencia distintos, enlazados únicamente por el morfismo tipado \(6\to8\). Las \(12\) aristas del cubo conservan su función combinatoria propia y no definen por sí solas la incidencia de Steiner. La diferencia de incidencia expresa una incompatibilidad dimensional productiva: la información de superficie requiere una regla de transducción para adquirir residencia volumétrica.

El mismo principio aparece en el acarreo \(999+1=1000\). Los nueve alcanzan el borde del registro decimal; el incremento publica simultáneamente el cierre del nivel anterior y el origen \(0/1\) del siguiente. La razón holográfica entre \(729\), \(243\), \(27\) y \(1\) despliega esa transición en escalas triádicas. La transducción de Barbero–Immirzi se inserta en esta familia de cambios de lector: frontera e interior conservan la unidad mientras redistribuyen su medida.

El defecto normalizado \(-1/12\) y el peso orientado \(12:-1\) ocupan funciones distintas dentro de esta geometría. El primero cuantifica una corrección de incidencia; el segundo registra una orientación firmada en la transducción. Su coordinación con el área mínima, el Hamiltoniano modular y la memoria proporciona el puente hacia la dinámica dimensional.

\subsection{Contracción hiperbólica interior, expansión exterior y conservación holográfica de la frontera}

La tesis puede formularse desde una idea única: la unidad conserva su identidad mientras cambia de escala, de lector, de curvatura, de memoria y de dimensión. Esta permanencia posee carácter evolutivo porque la unidad se despliega, adquiere nuevas coordenadas y publica propiedades diferentes; posee carácter holográfico porque cada publicación finita conserva la regla de reconstrucción del estado del que procede. El cierre holográfico del infinito designa precisamente esa compatibilidad entre identidad y despliegue: cada estadio contiene la ley que permite profundizar hacia dentro y prolongar hacia fuera la misma genealogía.

La denominación rectora es \textbf{principio de conservación evolutiva de la unidad mediante transducción dimensional y proyección holográfica}. La expresión sitúa cada término en su nivel correcto. La conservación pertenece a la identidad genealógica; la evolución pertenece a la actualización; la transducción relaciona realizaciones dimensionales; la proyección holográfica conserva la correspondencia reconstructiva entre sus lectores.

La aritmética elemental ofrece la primera imagen de esta arquitectura. La completación decimal

\[
1000=10^3=729+243+27+1=3^6+3^5+3^3+3^0
\]

expone cuatro estratos triádicos dentro de una unidad cúbica decimal. Al normalizar,

\[
1=\frac{729}{1000}+\frac{243}{1000}+\frac{27}{1000}+\frac{1}{1000},
\]

la unidad se reconoce simultáneamente como totalidad y como despliegue graduado. La lectura recíproca

\[
3^{-6},\quad 3^{-5},\quad 3^{-3},\quad 3^0
\]

describe la profundidad interior de esos estratos; su normalización produce la partición correspondiente. La relación \(999+1=1000\) expresa el mecanismo de transporte: la última unidad completa un registro y actúa, a la vez, como acarreo que inaugura el siguiente. El cero y el uno forman así una interfaz de continuidad. Los intervalos \([0,1]\), \([0,10]\) y la prolongación hacia \([0,\infty)\) son lectores escalares de una misma ley de completación.

La primera escisión \(1000=729+271\), seguida por \(271=243+27+1\), da forma a la razón extrema y media holográfica del despliegue: un núcleo triádico dominante conserva dentro del resto la misma regla de resolución. La unidad cúbica decimal contiene así una profundidad triádica anidada. El movimiento hacia \(729\), \(243\), \(27\) y \(1\) lee la contracción interior; el acarreo hacia el millar lee la expansión exterior.

Esta doble orientación organiza el infinito HMT. Hacia dentro, el refinamiento crece hiperbólicamente: cada celda admite una nueva profundidad sin perder su filiación. Hacia fuera, la unidad se expande mediante publicaciones dimensionales sucesivas. El cierre consiste en la conmutatividad de ambos movimientos. Si \(U\) representa la actualización del estado, \(R_d\) el lector en el estrato \(d\) y \(T_{d\to d'}\) la transducción entre estratos, la conservación evolutiva se expresa mediante

\[
R_{d'}\circ U
=
T_{d\to d'}\circ R_d.
\]

El mismo acontecimiento puede ofrecer una medida angular, una densidad de acción, una masa publicada o una respuesta gravitatoria. Su identidad común reside en la genealogía que hace conmutar esas lecturas.

\subsubsection{Monodromía nonádica y elevación helicoidal con avance de memoria}

APP, TRIT y TPK proporcionan el mecanismo formal de este despliegue. APP construye el soporte y separa las evaluaciones que deben conservarse. TRIT organiza los regímenes locales y su ambivalencia. TPK transporta fase y memoria sobre el soporte enriquecido. La dependencia causal queda fijada por

\[
U_t\longrightarrow \Gamma_9\longrightarrow K_{\mathrm{ph}}.
\]

El retorno de fase y el avance de memoria forman operaciones coordinadas. Tras nueve pasos puede reaparecer la misma clase angular,

\[
g(K+9)=g(K),
\]

mientras el contador de memoria progresa,

\[
q_{\mathrm{mem}}\Gamma_9=q_{\mathrm{mem}}+1,
\qquad
B_{K+9}\neq B_K.
\]

La figura geométrica precisa es una elevación helicoidal. Una órbita cerrada en la proyección angular se eleva a una trayectoria abierta cuando cada retorno de fase incrementa la memoria. La circunferencia conserva el ciclo; la hélice conserva el ciclo y registra además la profundidad alcanzada. La monodromía coinductiva genera así profundidad infinita mediante una sucesión de retornos localmente cerrados y globalmente ascendentes.

Esta imagen supera el valor meramente ilustrativo. La rueda angular constituye una sección de la hélice; la hélice constituye la biografía completa de la rueda cuando el estado transporta memoria. Cada vuelta recupera una posición de fase y, simultáneamente, inaugura una hoja nueva. La infinitud procede de la compatibilidad entre retorno y elevación: reiterar el ciclo produce profundidad sin disolver la identidad orbital.

El certificado nonádico da contenido finito y verificable a esa arquitectura. La cadena

\[
w_6\longrightarrow w_{12}\longrightarrow w_{18}\longrightarrow
w_{24}\longrightarrow w_{30}\longrightarrow R_{36}\longrightarrow G_9
\]

ordena el levantamiento; los \(396\) niveles, los \(2376\) trits, las \(43\) vueltas y los \(387\) pares \((K,K+9)\) por canal cuantifican el tránsito; las cinco regiones de \(\pi\) y la fibra ambiente de \(729\) elementos publican lecturas internas posteriores. Esta secuencia hace visible una monodromía productiva: el retorno recupera la fase y la coinducción conserva el crecimiento de memoria.

\subsubsection{Unidad radional y coordinación de los lectores circular, elíptico, proyectivo y helicoidal}

El radión concentra en una unidad metrológica la estructura que el radián publica angularmente. Su definición natural es la de una \textbf{sección orientada de fase, acción y memoria}. Una unidad de fase \(\theta=1\) transporta una unidad de acción \(\sigma\hbar_{\mathrm{ret}}\), pertenece a una vuelta completa \(T_+=2\pi_{\mathrm{HMT}}\) y conserva hoja, orientación, retorno nonádico y posición en el cilindro de refinamiento.

El radián expresa la coordenada angular visible. El radión conserva, junto a esa coordenada, la acción transportada, la orientación, el acarreo y la memoria de la vuelta. La metrología angular adquiere así una genealogía. Medir un arco significa identificar cuánto giro se ha publicado y qué estado completo ha realizado ese giro.

La relación entre fase y área adopta una forma especialmente clara:

\[
\mathcal A(\theta)=\hbar_{\mathrm{ret}}\,\theta.
\]

Una unidad radional barre un área de acción \(\hbar_{\mathrm{ret}}\); una vuelta completa encierra \(h_{\mathrm{ret}}\). El paso de fase a área queda incorporado en la propia unidad de lectura. La métrica angular, la acción y la memoria dejan de aparecer como cantidades yuxtapuestas y pasan a ser coordenadas de una misma sección orientada.

El radión admite varias publicaciones geométricas coordinadas. El lector circular publica fase; el lector elíptico publica acción y forma; el lector interior de Klein publica profundidad; el lector helicoidal nonádico publica retorno y memoria. La unidad permanece genealógicamente idéntica a través de esas realizaciones. Por eso el radión constituye el enlace metrológico adecuado entre HMT y Mecánica Dimensional: ofrece una unidad que puede atravesar cambios de lector conservando el estado que cada medida parcial presupone.

\subsection{Polaridad gravitatoria pentacomponente y transducción conjunta de acción, área, volumen, tiempo y masa}

La Mecánica Dimensional permite definir la dimensión como un régimen estable de publicación caracterizado por el rango mínimo de memoria independiente requerido por el levantamiento. Atravesar una dimensión significa aplicar un transductor que conserva la genealogía y cambia el tipo de lectura. La gravedad ocupa el nivel de conexión entre esos regímenes.

En esta formulación, la \textbf{gravitación interdimensional de cinco polarizaciones correlativas}, abreviadamente \textbf{estructura pentapolar de la gravitación}, es la conexión que transporta la unidad genealógica y deja en cada interfaz una torsión mínima. Una escritura compacta para esa traza es

\[
\Theta_d^{\min}
=
\nabla_{d+1}\circ T_{d\to d+1}
-
T_{d\to d+1}\circ\nabla_d.
\]

La expresión mide la diferencia entre transportar después de derivar y derivar después de transportar. Esa diferencia es la memoria geométrica mínima del cruce dimensional. La composición de interfaces produce una holonomía gravitatoria

\[
\mathfrak G_{d\to D}
=
\operatorname{Hol}
\bigl(
\Theta_d^{\min},
\Theta_{d+1}^{\min},
\ldots,
\Theta_{D-1}^{\min}
\bigr).
\]

El objeto fundamental es esta conexión interdimensional. Un modo local tipo gravitón ocupa el rango derivado de publicación cuantizada de la torsión; la gravitación conserva su residencia primaria en la compatibilidad entre dimensiones.

El término pentapolar recoge las cinco polarizaciones consustanciales del estado continuo:

\[
\mathbf G_d^{\pm}
=
\bigl(
G_{\mathrm{sp},d}^{\pm},
G_{\mathrm{sol},d}^{\pm},
G_{\mathrm{gau},d}^{\pm},
G_{\mathrm{coh},d}^{\pm},
G_{\mathrm{det},d}^{\pm}
\bigr).
\]

Las cinco componentes emergen correlativamente del mismo estado enriquecido. La polaridad \(\pm\) registra concentración directa y publicación conjugada. La gravedad atraviesa los estratos dimensionales transportando simultáneamente esas cinco lecturas; la torsión mínima conserva la huella de cada cruce.

La geometría adquiere con ello una imagen unificada. El régimen elíptico organiza cierres; el parabólico organiza umbrales; el hiperbólico organiza profundidad. La transducción dimensional enlaza esos regímenes con las respuestas radiales \(p\), las holonomías de espín y las publicaciones de masa. La estructura pentapolar de la gravitación expresa la coherencia global de ese enlace.

La teoría M proporciona pantallas físicas posteriores para esta conexión: dimensiones adicionales, cuerdas, branas y dualidad T realizan modos concretos de transducción. HMT aporta la genealogía APP–TRIT–TPK, el estado enriquecido, la memoria y la lectura bidireccional; la interfaz con teoría M publica esas estructuras en geometrías físicas tipadas. La dualidad T intercambia lectores dimensionales compatibles y conserva la clase genealógica del estado.

\subsubsection{Compatibilidad de las realizaciones clásica, relativista, cuántica y HMT--MD}

La mecánica clásica publica trayectoria y momento. La relatividad publica estructura causal, métrica y curvatura. La mecánica cuántica publica fase, holonomía, vacancias y amplitudes. La Mecánica Dimensional proporciona los transductores que relacionan esas publicaciones, y HMT conserva la genealogía discreta que las sostiene.

En la lectura clásica, una órbita describe la proyección visible del movimiento. En la lectura relativista, esa órbita queda incorporada a la geometría del lector. En la lectura cuántica, la fase y el espín registran la holonomía de la sección. En la lectura HMT–MD, las tres expresiones pertenecen a una misma biografía elevada: el radión cuantifica fase, acción y memoria; \(p\) clasifica la respuesta radial; \(\tau\) clasifica la curvatura; la holonomía publica espín; el transductor dimensional publica masa; la conexión pentapolar publica gravitación.

La aparente disipación clásica adquiere residencia en el balance genealógico ampliado. La aceleración relativista se expresa mediante conexión y proyección. La indeterminación cuántica se organiza por lectores, vacancias y ramas de Ambivalencia. El principio de Mach relaciona la respuesta local con la distribución global. El Landauer nulo garantiza la reversibilidad informacional del estado completo. Todas estas piezas convergen en una conservación más profunda: la identidad de la unidad a través de cambios de publicación.

\subsubsection{Teorema narrativo de conservación genealógica y prolongación holográfica}

El cierre puede narrarse ahora como una sola secuencia. La unidad se descompone sin perder identidad. La fase orientada adquiere medida en el radión. El retorno nonádico convierte el ciclo en hélice mediante memoria. El carácter \(p\) transforma la hélice en familias radiales y clasifica sus regímenes. La persistencia de esas trayectorias produce secciones con espín, carga, sabor, color y masa. Las vacancias abren canales; la polaridad eléctrica transporta fase; la respuesta magnética curva transversalmente; los umbrales publican radiación. La autoconservación cierra el balance genealógico de acción y momento. La autoinercia relacional enlaza conexión propia, proyección y totalidad machiana. La Ambivalencia conserva las direcciones directa y conjugada. La transducción área–volumen conduce a Barbero–Immirzi y a la ley de masas. La torsión mínima encadena dimensiones y publica la estructura pentapolar de la gravitación.

El crecimiento hacia dentro y la expansión hacia fuera son orientaciones conjugadas del mismo proceso. El refinamiento hiperbólico genera profundidad; la transducción dimensional genera extensión. El acarreo \(0/1\) enlaza escalas; la monodromía enlaza vueltas; la memoria enlaza estados; la holonomía enlaza orientaciones; la gravitación enlaza dimensiones. Cada enlace conserva una unidad que evoluciona y cada publicación finita contiene la regla de su prolongación.

La tesis fundamental queda formulada así:

\begin{quote}
\textbf{El cierre holográfico del infinito realiza el principio de conservación evolutiva de la unidad mediante transducción dimensional y proyección holográfica: la unidad mantiene su identidad genealógica al desplegarse como fase, acción, curvatura, memoria, partícula, campo y gravitación; cada lectura finita conserva la regla de reconstrucción y prolongación del estado completo.}
\end{quote}

Esta formulación ofrece una imagen física y matemática conjunta. El infinito adquiere una ley de cierre local; la unidad adquiere profundidad; la dimensión adquiere memoria; la conservación adquiere carácter relacional; la fuerza adquiere autoconservación; la inercia adquiere estructura machiana; la gravedad adquiere conexión pentapolar; y la holografía adquiere el sentido preciso de reconstrucción entre lectores compatibles.


\section{Definición operatoria del radión angular mediante el estado enriquecido y \(4\) lectores tipados}
\subsection{Definición operatoria del radión angular y cuádruple realización tipada del estado enriquecido}
% Fragmento público generado desde propietarios íntegros.
% Residencia: 52.13.
% Fuente: ../incoming/radion_monografia_20260827/manuscrito/sections/13_sintesis.tex:1-80
\noindent La síntesis operatoria reúne los lectores de fase, acción, profundidad y memoria sin identificarlos entre sí.\par\medskip

\subsubsection{Cuádruple realización del estado radional}

Sea \(x_\infty\in X_{\mathrm{enr}}\) la sección coinductiva. El radión completo no
es un punto aislado, sino el paquete
\begin{equation}
\mathfrak R_\sigma(x_\infty)
=\bigl(
\theta=1,\sigma,
\pih,\alphah,A,C^*,\Delta_4,
\Phi_T,D_E,\epsone,\hret,
q_{\mathrm{mem}},\mathcal I_\infty
\bigr).
\label{eq:full-radion}
\end{equation}
Cuatro lectores extraen de él objetos diferentes:

\begin{figure}[H]
\centering
\begin{tikzpicture}[
  node distance=17mm and 18mm,
  state/.style={draw=RadNavy,very thick,rounded corners,fill=RadCream,
    minimum width=39mm,minimum height=13mm,align=center},
  readout/.style={draw=RadTeal,rounded corners,fill=RadMist,
    minimum width=39mm,minimum height=12mm,align=center},
  arr/.style={-{Latex},thick,RadCopper}]
\node[state] (x) {estado enriquecido\\\(x_\infty\)};
\node[readout,above left=of x] (circ) {círculo \(S^1\)\\fase};
\node[readout,above right=of x] (ell) {elipse positiva\\acción y forma};
\node[readout,below left=of x] (klein) {interior Klein\\profundidad};
\node[readout,below right=of x] (helix) {hélice nonádica\\memoria y retorno};
\draw[arr] (x)--(circ) node[midway,below left]{\small \(P_{\mathrm{ph}}\)};
\draw[arr] (x)--(ell) node[midway,below right]{\small \(P_{\mathrm{act}}\)};
\draw[arr] (x)--(klein) node[midway,above left]{\small \(P_K\)};
\draw[arr] (x)--(helix) node[midway,above right]{\small \(P_{\mathrm{mem}}\)};
\end{tikzpicture}
\caption{Cuatro publicaciones tipadas del mismo antecedente. La unidad del
estado no elimina la diferencia de codominios.}
\label{fig:four-readers}
\end{figure}

El círculo conserva la clase de fase y olvida el lift. La elipse conserva
acción, anisotropía y orientación simpléctica. Klein dota al interior de una
distancia proyectiva. La hélice conserva número de retornos y refinamiento.
La Mecánica Dimensional no consiste en mezclarlos, sino en construir los
transductores que los hacen compatibles.

\subsubsection{Interpretación operatoria del cociclo \(\varepsilon_R\)}

\(\varepsilon_R\) no es una constante universal independiente de todo contexto. Es la
coordenada de un cociclo de transición asociado a la carta del radión. Su
dominio contiene dos secciones del mismo estado:
\[
S_E=1+D_E,
\qquad
S_T=1+\Phi_T.
\]
Su regla es
\begin{equation}
\epsone=S_T-S_E=\Phi_T-D_E.
\label{eq:epsilon-final-meaning}
\end{equation}
Su codominio es la recta adimensional de unidad angular; la carta
\(360/T_+\) la publica en grados.

\begin{exactbox}
\textbf{En lenguaje llano:} la sección directa y la sección torsional llegan
a la misma celda de unidad por dos rutas internas diferentes. Sus partes
enteras coinciden; sus restos no. \(\epsR\) es la diferencia orientada de esos
restos, calculada con el acarreo de APP y llevada al límite por los cilindros
TPK. No se introduce para que el radián cierre: aparece antes de consultar el
radián y, cuando se publica, lo cierra.
\end{exactbox}

Esta definición explica por qué el número es pequeño. No porque se haya
buscado una corrección pequeña, sino porque \(S_E\) y \(S_T\) comparten un
prefijo profundo. La pequeñez mide la cercanía estructural de dos
publicaciones; el signo mide su orientación; las cifras registran el acarreo.


% Fragmento público generado desde propietarios íntegros.
% Residencia: 52.13.1.
% Fuente: ../incoming/radion_monografia_20260827/manuscrito/sections/13_sintesis.tex:81-134
\subsubsection{Definición operatoria del radión}

\begin{definition}[Radión HMT--MD]
El radión es la sección orientada de fase, acción y memoria que cumple
simultáneamente:
\begin{enumerate}
\item recorre una unidad de fase \(\theta=1\);
\item barre acción firmada \(\sigma\hret\);
\item pertenece a una vuelta \(T_+=2\pih\) generada coinductivamente;
\item conserva la hoja \(\sigma\), el retorno nonádico y el cilindro de
refinamiento;
\item satisface el empalme exacto
\(1=S_E-\Phi_T+\epsone\).
\end{enumerate}
\end{definition}

El radián convencional es la imagen del radión bajo el functor de olvido
\[
\mathfrak R_\sigma
\longmapsto
\theta=1\in\R/2\pi\Z.
\]
Ese olvido es útil y legítimo para la trigonometría euclídea. Se vuelve
insuficiente cuando la física necesita distinguir acción, signo, retorno,
campo e historia.

\paragraph{Genealogía de la unidad angular natural}

La naturalidad del radián suele justificarse porque la derivada de
\(e^{i\theta}\), \(\sin\theta\) y \(\cos\theta\) adopta su forma más simple
cuando \(\theta\) se mide como cociente arco/radio. HMT conserva esa ventaja,
pero la explica desde una capa anterior:
\[
\theta\quad\leftrightarrow\quad
\frac{\mathcal A(\theta)}{\hret}.
\]
La fase es natural porque mide acción en unidades de \(\hret\); la vuelta es
natural porque encierra \(\hfull\). El cociente arco/radio es la publicación
euclídea de una unidad que ya posee estructura simpléctica y memoria.

\paragraph{Compatibilidad entre elipse de acción y crecimiento dimensional}

El círculo es la forma que resulta al olvidar la anisotropía:
\(C^*=0\Rightarrow\eta=0\Rightarrow a=b\). La elipse es más informativa porque
conserva dos autovalores, dos ejes y una orientación de intercambio a producto
fijo. En ese sentido, la elipse es una prolongación del círculo, no un círculo
deformado por accidente.

La profundidad de Klein y el lift helicoidal añaden dos tipos de información
que el contorno plano no contiene: distancia interior y memoria de retorno.
Ésta es la imagen ontológica de Mecánica Dimensional que emerge del radión:
una dimensión nueva aparece cuando se exige conservar un dato independiente
que la proyección anterior debía olvidar.


\subsection{Propiedades estructurales de la unidad radional: genealogía, bisagra crítica, densidad, acarreo, acción, compatibilidad borde--interior, forma, polaridad y modularidad}
% Fragmento público generado desde propietarios íntegros.
% Residencia: 52.13.2.
% Fuente: ../incoming/radion_monografia_20260827/manuscrito/sections/13_sintesis.tex:135-184
\subsubsection{Sistema de teoremas componentes del radión}

\begin{theorem}[Genealogía]
El objeto \(\mathfrak R_\sigma\) desciende de
\(\APP\to\TRIT\to\TPK\to X_{\mathrm{enr}}\) y no recibe el blanco metrológico
como selector.
\end{theorem}

\begin{theorem}[Bisagra]
En la carta APP de cien marcas, la costura crítica y la fase dodecafásica
segunda satisfacen \(j_{100}(1/2)=\theta_2=50^\circ\), con \(m=2\) único.
\end{theorem}

\begin{theorem}[Densidad]
Las dos hojas del sector activo \(N_6=90\), normalizadas por \(3^6=729\),
producen \(\rho_{\mathrm{tw}}=20/81\).
\end{theorem}

\begin{theorem}[Acarreo]
La resta APP de las secciones \(S_T\) y \(S_E\) converge a
\(\epsone\) con tasa de intervalo \(3^{-54}\) por retorno y cierra exactamente
la unidad.
\end{theorem}

\begin{theorem}[Acción]
La elipse positiva satisface \(\mathcal A(\theta)=\hret\theta\); un radión
barre \(\hret\) y una vuelta encierra \(\hfull\).
\end{theorem}

\begin{theorem}[Borde--interior]
La misma frontera posee acción finita \(\hfull\), mientras el interior de
Klein tiene área hiperbólica infinita.
\end{theorem}

\begin{theorem}[Forma]
El cociente \(e^{-\eta}\) se conserva a través de semiejes espectrales,
semiejes de acción, incertidumbres, impedancia y módulo del toro.
\end{theorem}

\begin{theorem}[Polaridad]
La inversión \(C^*\mapsto-C^*\) intercambia constitutivas, invierte
impedancia y conserva la velocidad normalizada.
\end{theorem}

\begin{theorem}[Modularidad]
La inversión bidireccional actúa como \(S:\tau\mapsto-1/\tau\) y conserva
\(j(\tau)\), mientras el etiquetado focal conserva la orientación que
\(j\) olvida.
\end{theorem}

% Fuente: ../incoming/radion_monografia_20260827/manuscrito/sections/A_demostraciones.tex:1-179
\subsubsection{Demostraciones algebraicas y geométricas reunidas}

\paragraph{Demostración lineal del cierre unitario}

Para facilitar una auditoría independiente, reunimos el argumento sin
interrupción narrativa.

\begin{enumerate}[label=\textbf{Paso \arabic*.},leftmargin=2.2cm]
\item La conexión nonádica genera correlacionadamente
\(\pih,e_{\HMT},\varphi_{\HMT},\alphah\).
\item La vuelta es \(T_+=2\pih\).
\item La rueda fija \(\theta_2=50^\circ\), y la carta parangular fija
\(A=1000\alphah\).
\item La sección directa es
\[
S_E=T_+\left(\frac{50}{360}+\frac{1000\alphah}{360}\right)
=2\pih\left(\frac5{36}+\frac{25}{9}\alphah\right).
\]
\item Los intervalos certificados dan \(1<S_E<2\); luego
\(D_E=S_E-1\).
\item La incidencia produce \(N_6=90\); dos hojas y el ambiente \(729\)
producen \(\rho=2N_6/729=20/81\).
\item La torsión global es
\[
\Delta_4=\frac{\pih-e_{\HMT}\log\pih}{270}
\left(1+\frac{\pih}{729}\right).
\]
\item La sección torsional es \(S_T=1+(20/81)\Delta_4\).
\item La resta APP define
\[
\epsone=S_T-S_E=\frac{20}{81}\Delta_4-D_E.
\]
\item Por sustitución,
\[
S_E-\frac{20}{81}\Delta_4+\epsone=1.
\]
\item Multiplicar por \(360/T_+=180/\pih\) produce
\[
\frac{180^\circ}{\pih}
=50^\circ+A^\circ-\frac{20}{81}\Delta_4^\circ+\epsR^\circ.
\]
\end{enumerate}

Ningún paso recibe el lado izquierdo de la última igualdad. El resultado
metrológico sólo se evalúa al final como reconocimiento.

\paragraph{Unicidad de la resta con acarreo}

Sea \(B=1000\). Para cada posición, todo entero \(u_i\) admite una división
euclídea única
\[
u_i=Bc_i+r_i,\qquad 0\le r_i<B.
\]
Por tanto,
\[
r_i=u_i\bmod B,\qquad c_i=(u_i-r_i)/B
\]
son únicos. Recorriendo las posiciones desde el extremo remoto hacia el
frente, el acarreo de salida de una posición fija la entrada de la siguiente.
Por inducción, toda la palabra de restos y carries es única.

En pseudocódigo:
\begin{verbatim}
acarreo = remote_acarreo
for i from last_block downto first_block:
    u = torsion[i] - direct[i] + acarreo
    remainder[i] = u mod 1000
    acarreo = (u - remainder[i]) / 1000
return remainder, acarreo_word
\end{verbatim}
El procedimiento no consulta el valor esperado de la diferencia.

\paragraph{Convergencia del funcional por profundidad}

Sea
\[
F(p,e)=1+\frac{20}{81}
\frac{p-e\log p}{270}\left(1+\frac{p}{729}\right)
-2p\left(\frac5{36}+\frac{25}{9}\alphah\right).
\]
En el rectángulo compacto \(3\le p\le4\), \(2\le e\le3\), \(F\) es
continuo y continuamente diferenciable. Si \(I_N^\pi\searrow\{\pih\}\) e
\(I_N^e\searrow\{e_{\HMT}\}\), la imagen intervalar
\[
F(I_N^\pi,I_N^e)
\]
forma una familia encajada cuya anchura tiende a cero. La intersección es un
único real, \(\epsone\). La monotonía de las derivadas permite evaluar los
extremos sin muestreo interior.

La tasa \(3^{-54}\) por retorno deriva del refinamiento de seis trits a lo
largo de nueve niveles. A profundidad \(45\) tríadas, la cota de anchura
\(<4.81\times10^{-130}\) certifica más de 129 cifras del valor.

\paragraph{Derivación de la métrica y del área hiperbólica de Klein}

En el disco unitario, la forma matricial de la métrica Klein es
\[
g(u,v)=\frac1{1-r^2}I
+\frac1{(1-r^2)^2}
\begin{pmatrix}u\\v\end{pmatrix}
\begin{pmatrix}u&v\end{pmatrix}.
\]
El lema del determinante para una perturbación de rango uno da
\[
\det g=(1-r^2)^{-3}.
\]
De ahí
\[
\sqrt{\det g}\,\diff u\diff v
=(1-r^2)^{-3/2}\diff u\diff v.
\]
Integrar en polares produce \eqref{eq:klein-area-R}. El límite infinito se
debe a la singularidad métrica en la frontera; la frontera sigue teniendo
área simpléctica finita porque esa medida utiliza \(\diff Q\wedge\diff P\),
no \(\sqrt{\det g}\).

\paragraph{Distancia focal en el disco de Klein}

En un diámetro, la distancia del centro a \(r\) es
\[
d_K(0,r)=\operatorname{artanh}r.
\]
Para \(r=e_f\), \eqref{eq:excentricidad} implica
\[
1-e_f^2=e^{-2\eta}.
\]
Si \(u_f=\operatorname{artanh}e_f\), entonces
\[
\operatorname{sech}u_f=\sqrt{1-e_f^2}=e^{-\eta},
\]
y por tanto \(\eta=\log\cosh u_f\), como en
\eqref{eq:focal-hyperbolic-chain}.

\paragraph{Transformación simpléctica de compresión}

Definamos
\begin{equation}
S_\eta=
\begin{pmatrix}e^{\eta/2}&0\\0&e^{-\eta/2}\end{pmatrix},
\qquad \det S_\eta=1.
\label{eq:compresión simpléctica}
\end{equation}
La estructura compleja transportada es
\begin{equation}
J_\eta=S_\eta
\begin{pmatrix}0&-1\\1&0\end{pmatrix}
S_\eta^{-1}
=\begin{pmatrix}0&-e^\eta\\e^{-\eta}&0\end{pmatrix},
\qquad J_\eta^2=-I.
\label{eq:Jeta}
\end{equation}
La curva
\[
\gamma(\theta)=\sqrt{2\hret}\,
S_\eta(\cos\theta,\sin\theta)
=e^{\theta J_\eta}\gamma(0)
\]
es exactamente la elipse de acción. El compresión simpléctica hiperbólico conserva área y
transporta la rotación compleja en vez de destruirla.

Para orientación \(\sigma\),
\[
\gamma_\sigma(\theta)
=(a_S\cos\theta,\sigma b_S\sin\theta),
\qquad
\mathcal A_\sigma(\theta)=\sigma\hret\theta.
\]
Las dos hojas comparten focos y módulo de acción; difieren en acción firmada.

\paragraph{Comprobación diofántica del peso dodecafásico de Barbero--Immirzi}

La ecuación \(4a-3b=45\) tiene soluciones enteras
\[
(a,b)=(12+3t,1+4t),\qquad t\in\Z.
\]
La condición \(|b|=1\) deja \(b=1\), pues \(b=-1\) no pertenece a la clase
\(1\pmod4\). De ahí \(a=12\). Esta comprobación reproduce el par obtenido mediante la cadena fundamental
dodecafásica: doce sectores y una única costura orientada. Los cardinales
\(90/120\), las multiplicidades de la cadena y el lector de retorno
conservan funciones distintas; el coeficiente de polo y la minimalidad
verifican aritméticamente su composición.

\subsection{Sensibilidad del cierre unitario frente a variaciones controladas de la multiplicidad de hojas, la fase dodecafásica, el canal \(e\), la torsión global, el operador armónico y la cola sexádica}
% Fragmento público generado desde propietarios íntegros.
% Residencia: 52.13.3.
% Fuente: ../incoming/radion_monografia_20260827/manuscrito/sections/13_sintesis.tex:185-247
\subsubsection{Controles negativos del sistema radional}

La construcción es refutable en puntos concretos:
\begin{enumerate}
\item Si \(S_E\notin(1,2)\), el cociente APP no es uno y el cierre cambia de
celda.
\item Si la incidencia activa no produce \(N_6=90\), o si no hay dos hojas,
\(20/81\) deja de seguirse.
\item Si el verificador recibe \(180/\pi\) o \(\epsR\) antes de construir los
operandos, la generación es circular.
\item Si la resta por tríadas no estabiliza prefijos compatibles, no existe
la sección coinductiva propuesta.
\item Si \(|C^*|\ge A\), la elipse positiva deja de existir.
\item Si \(a_Sb_S\ne2\hret\), falla el teorema de área.
\item Si \(C^*=0\), deben cumplirse \(d=0\), \(\eta=0\), \(\tau=i\) y
\(j=1728\).
\item Si el área Klein del subdisco no sigue \eqref{eq:klein-area-R}, falla la
tesis de profundidad infinita.
\item Si \(F_{\mathrm{spec}}=\epsR^\circ\) se afirma sin contratérmino independiente, la
diferencia numérica \eqref{eq:chiR} lo refuta.
\item La producción de \(12:-1\) exige conservar, junto a los cardinales
\(90/120\), los doce sectores dodecafásicos, la única costura orientada y
la acción del lector sobre esos generadores. Cambiar alguna multiplicidad
debe modificar la imagen \(12S_{90}-S_{120}\). Para la cadena fundamental
sin alteración, el coeficiente de \((1-q)^{-1}\) debe resultar \(1/24\).
\item Si la función \(j\) se usa para seleccionar signo de carga, se pierde
la invariancia \eqref{eq:j-invariance} y se incurre en un error de tipo.
\end{enumerate}

\subsubsection{Realización física conjunta del radión}

La imagen final no es un círculo corregido. Es una arquitectura de
publicaciones:
\[
\begin{aligned}
&\text{APP:}&&\text{unidad, resto, acarreo y fibra};\\
&\text{TRIT:}&&\text{régimen, orientación y hojas};\\
&\text{TPK:}&&\text{fase, retorno, dodecafase y memoria};\\
&\text{elipse:}&&\text{acción y anisotropía};\\
&\text{Klein:}&&\text{profundidad interior};\\
&\text{hélice:}&&\text{biografía del retorno};\\
&\text{MD:}&&\text{reloj, energía, carga, campo y causalidad}.
\end{aligned}
\]

La ecuación del radión reúne la cara visible:
\BigRadionEquation
La elipse reúne la cara de acción:
\[
\mathcal A(1)=\hret,\qquad \mathcal A(2\pi)=\hfull.
\]
La métrica de Klein reúne la profundidad:
\[
K=-1,\qquad \operatorname{Area}_K=\infty.
\]
La holonomía reúne la memoria:
\[
g(k+9)=g(k),\qquad B_{k+9}\ne B_k.
\]

\begin{thesisbox}
El círculo es la sombra angular; la elipse es el cuerpo de acción; Klein mide
el interior; la hélice conserva la biografía. El radión es la unidad que
permite leer las cuatro afirmaciones sobre un único estado sin reducir una a
otra.
\end{thesisbox}
% Fuente: ../incoming/radion_monografia_20260827/manuscrito/sections/B_valores_y_verificacion.tex:1-105
\subsubsection{Valores canónicos, estabilidad y verificación reproducible}

\paragraph{Valores de los invariantes internos generados}

\begin{longtable}{p{.27\textwidth} p{.61\textwidth}}
\toprule
Objeto & Valor usado en el perfil coinductivo\\
\midrule
\endfirsthead
\toprule
Objeto & Valor usado en el perfil coinductivo\\
\midrule
\endhead
\(\pih\) & \(\begin{aligned}3.1415926535897932384626433832795028841\\[-2pt]971693993751058209749445923\ldots\end{aligned}\)\\
\(e_{\HMT}\) & evaluación coinductiva del lector de Euler, sin tabla externa\\
\(\alphah\) & \(\begin{aligned}0.0072973525692838009972851054723806629\\[-2pt]995641725396427091854046825\ldots\end{aligned}\)\\
\(A\) & \(\begin{aligned}7.2973525692838009972851054723806629995\\[-2pt]641725396427\ldots{}^\circ\end{aligned}\)\\
\(C^*\) & \(2.1237383389686457242619513803933607937\ldots{}^\circ\)\\
\(\eta\) & \(0.299689683921630799684926562863927\ldots\)\\
\(\hret\) & \(1.05457181691503906113369058472430\ldots\times10^{-34}U_S\)\\
\bottomrule
\end{longtable}

\paragraph{Convergencia cuantitativa por profundidad}

\begin{center}
\small
\begin{tabular}{r l l}
\toprule
Profundidad & \(\epsR^\circ\) & lectura\\
\midrule
12 & \(5.1681282186551375597\times10^{-8}\) & preasintótica\\
18 & \(5.1383017764316644535\times10^{-8}\) & 7 cifras estables\\
24 & \(5.1383016758575394560\times10^{-8}\) & 14 cifras estables\\
30 & \(5.1383016758575282072\times10^{-8}\) & 20 cifras estables\\
36 & \(5.138301675857528207168517\times10^{-8}\) & perfil largo\\
45 & \(5.13830167585752820716851646226459\times10^{-8}\) & ancho \(<4.81\times10^{-130}\)\\
\bottomrule
\end{tabular}
\end{center}

Los retornos de profundidad \(9,18,27,36,45\) conservan la fase \(8\) y
aumentan la memoria \(0,1,2,3,4\). La razón entre anchos consecutivos tiende
a
\[
3^{-54}=1.7196982334549856127610399316004871\ldots\times10^{-26}.
\]

\paragraph{Análisis de sensibilidad estructural por supresión controlada}

\begin{longtable}{p{.46\textwidth} p{.34\textwidth}}
\toprule
Operación & Salida o diferencia unitaria\\
\midrule
\endfirsthead
\toprule
Operación & Salida o diferencia unitaria\\
\midrule
\endhead
Operador canónico & \(8.96802822044562981999\times10^{-10}\)\\
Una hoja, \(\rho=90/729\) & \(-1.3727056615960678\times10^{-5}\)\\
Fase \(\theta_1=20^\circ\) & \(+5.2359877649510169\times10^{-1}\)\\
Fase \(\theta_3=80^\circ\) & \(-5.2359877470149605\times10^{-1}\)\\
Sin canal \(e\) & \(1.8065350748904653\times10^{-3}\)\\
Torsión histórica \(\Delta_4(A,C^*)\) & \(1.0517084608768816\times10^{-9}\)\\
Operador armónico \(M_4\) menos canónico & \(1.0525500222895070\times10^{-17}\)\\
Cola sexádica menos canónico & \(2.0283936357433839\times10^{-12}\)\\
\bottomrule
\end{longtable}

La separación de escalas muestra sensibilidad estructural: el cierre depende
de la fase, las dos hojas, los canales \(\pi/e\), la torsión global y el acarreo.
Ninguna ablación reproduce el valor.

\paragraph{Certificados computacionales reproducibles}

La obra usa los siguientes propietarios ejecutables:
\begin{enumerate}
\item \sourcepath{output/RADION_TWOWAY_NONADICO_20260827/verificar_emergencia_radion_por_profundidad.py};
\item \sourcepath{output/RADION_TWOWAY_NONADICO_20260827/verificar_radion_helice_critica_tpk.py};
\item \sourcepath{output/RADION_TWOWAY_NONADICO_20260827/verificar_operador_tpk_target_free_profundidad.py};
\item el verificador consolidado de esta monografía, ubicado en
\sourcepath{certificados/verificar_monografia_radion.py}.
\end{enumerate}
Las salidas esperadas incluyen
\[
\texttt{PASS\_EMERGENCIA\_RADION\_POR\_PROFUNDIDAD},
\]
\[
\texttt{PASS\_RADION\_HELICE\_CRITICA\_TPK},
\qquad
\texttt{PASS\_MONOGRAFIA\_RADION\_HMT\_MD}.
\]

\paragraph{Distinción entre precisión computacional y exactitud matemática}

Que un programa muestre cero a 120 o 200 cifras no convierte por sí solo una
identidad residual en predicción. La exactitud matemática procede de
\eqref{eq:cierre-unitario}; la precisión arbitraria permite evaluar sus
operandos y verificar la convergencia. La no circularidad procede del orden
de construcción y de la prohibición de blancos, no del número de dígitos.

El cálculo Decimal puede dejar residuos del orden \(10^{-218}\) al combinar
ramas evaluadas con precisiones distintas. Ese número es redondeo de la
implementación; no una corrección física adicional.
% Fuente: ../incoming/radion_monografia_20260827/manuscrito/sections/D_historia_fuentes.tex:1-169
\subsubsection{Genealogía conceptual y fuentes primarias}

\paragraph{Cuanto de acción de Planck}

Planck introduce \(h\) en la ley de radiación como constante que liga energía
y frecuencia. El dato histórico decisivo para esta obra es dimensional:
\(h\) es acción. El radión no modifica ese papel; explica por qué la acción de
una vuelta y la acción por unidad de fase se relacionan mediante \(2\pi\).

\paragraph{Cuantización angular de Bohr y Sommerfeld}

Bohr usa \(h/2\pi\) en la cuantización del momento angular. Sommerfeld amplía
la cuantización a integrales de acción y órbitas keplerianas. Por ello, no es
históricamente exacto decir que Dirac inventó la constante reducida. Su
contribución pertenece a una genealogía en la que la acción angular ya estaba
presente.

La novedad HMT--MD es distinta: no toma \(h/2\pi\) como un cociente externo,
sino que construye una elipse cuya vuelta tiene área \(h\) y cuyo sector de un
radián tiene área \(\hbar\).

\paragraph{Formalismo matricial de Born y Jordan y oscilador cuántico}

Born y Jordan sitúan el principio variacional, las fases y el oscilador
armónico en el corazón de la mecánica matricial; el campo electromagnético se
descompone en una colección de osciladores. La realización LC del radión
continúa esa intuición con un dato añadido: el cociente de forma
\(e^{-\eta}\) determina la impedancia relativa mientras producto, frecuencia
y acción permanecen fijos.

\paragraph{Frecuencia angular y acción en la formulación de Dirac}

En 1925, Dirac escribe las frecuencias en radianes por unidad de tiempo y usa
relaciones de la forma
\[
xy-yx=\frac{ih}{2\pi}[x,y],
\qquad
\Delta E=\frac{h}{2\pi}\omega.
\]
La ecuación \(E=\hret\omega\) une de manera explícita acción reducida y
frecuencia angular. En el radión, esa relación recibe una geometría: \(\hret\)
es el área barrida por una unidad de fase.

\paragraph{Relaciones de indeterminación de Heisenberg, Robertson y Schrödinger}

Heisenberg formula la limitación conjunta de variables canónicamente
conjugadas. Robertson generaliza la desigualdad; Schrödinger incorpora la
covarianza. La matriz \eqref{eq:covariance} pertenece más directamente a esa
prolongación Robertson--Schrödinger que al artículo de Heisenberg aislado.

Ninguno de estos autores formula la elipse genealógica completa de esta obra.
La afirmación históricamente defendible es que proporcionan el lenguaje en el
que la elipse de acción puede reconocerse como un contorno mínimo; la
construcción del par \((A,C^*)\), los focos, Klein, la hélice y el acarreo es la
aportación HMT--MD.

\paragraph{Funcional de fase de Feynman}

La formulación espacio-temporal asigna a cada camino una fase
\(e^{iS/\hbar}\). La unidad de radión hace literal la correspondencia:
incrementar la acción en \(\hbar\) incrementa la fase en un radián. La suma de
caminos reconoce, en un lenguaje integral, el mismo principio que la elipse
expresa como área.

\paragraph{Propagaciones avanzada y retardada en Wheeler--Feynman}

La teoría del absorbedor construye una interacción clásica temporalmente
simétrica mediante combinaciones de campos retardados y avanzados. La
identidad HMT \(CU(t)C^{-1}=U(-t)\) y la inversión de rutas proporcionan un
antecedente operatorio. La promoción física requiere que el mapa entre rutas
y propagadores preserve acción, orientación, frontera y forma simpléctica.

La teoría del absorbedor no es la misma afirmación que la interpretación de
Feynman y Stückelberg de las antipartículas como propagación temporalmente
invertida. Ambas pueden reconocer la bidireccionalidad; sus objetos físicos y
condiciones de frontera son diferentes.

\paragraph{Dualidad de Hodge y separación respecto de la cohomología de Hochschild}

La operación relevante en este manuscrito es la estrella de Hodge
\(\star^2=-I\) sobre dos-formas en firma lorentziana. No se ha localizado en
el corpus activo del radión un propietario de homología de Hochschild que
actúe sobre la elipse, el acarreo o el sector \(90/120\). Introducir Hochschild
por semejanza nominal degradaría el tipado. Si una futura prolongación
construye un cociclo de deformación con dominio y codominio propios, deberá
entrar como resultado nuevo, no como retroatribución.

\paragraph{Bibliografía primaria seleccionada}

\begin{enumerate}
\renewcommand{\labelenumi}{[\arabic{enumi}]}
\item\hypertarget{radionbib:Planck1901}{}
M.~Planck,
\emph{Ueber das Gesetz der Energieverteilung im Normalspectrum},
Annalen der Physik 309 (1901), 553--563.
\url{https://doi.org/10.1002/andp.19013090310}

\item\hypertarget{radionbib:Bohr1913}{}
N.~Bohr,
\emph{On the Constitution of Atoms and Molecules},
Philosophical Magazine 26 (1913), 1--25.
\url{https://doi.org/10.1080/14786441308634955}

\item\hypertarget{radionbib:Sommerfeld1916}{}
A.~Sommerfeld,
\emph{Zur Quantentheorie der Spektrallinien, I},
Annalen der Physik 356 (1916), 1--94.
\url{https://doi.org/10.1002/andp.19163561702}

\item\hypertarget{radionbib:BornJordan1925}{}
M.~Born y P.~Jordan,
\emph{Zur Quantenmechanik},
Zeitschrift für Physik 34 (1925), 858--888.
\url{https://doi.org/10.1007/BF01328531}

\item\hypertarget{radionbib:Dirac1925}{}
P.~A.~M.~Dirac,
\emph{The Fundamental Equations of Quantum Mechanics},
Proceedings of the Royal Society A 109 (1925), 642--653.
\url{https://doi.org/10.1098/rspa.1925.0150}

\item\hypertarget{radionbib:Heisenberg1927}{}
W.~Heisenberg,
\emph{Über den anschaulichen Inhalt der quantentheoretischen Kinematik und Mechanik},
Zeitschrift für Physik 43 (1927), 172--198.
\url{https://doi.org/10.1007/BF01397280}

\item\hypertarget{radionbib:Dirac1928}{}
P.~A.~M.~Dirac,
\emph{The Quantum Theory of the Electron},
Proceedings of the Royal Society A 117 (1928), 610--624.
\url{https://doi.org/10.1098/rspa.1928.0023}

\item\hypertarget{radionbib:Robertson1929}{}
H.~P.~Robertson,
\emph{The Uncertainty Principle},
Physical Review 34 (1929), 163--164.
\url{https://doi.org/10.1103/PhysRev.34.163}

\item\hypertarget{radionbib:Schrodinger1930}{}
E.~Schrödinger,
\emph{Zum Heisenbergschen Unschärfeprinzip},
Sitzungsberichte der Preussischen Akademie der Wissenschaften (1930), 296--303.

\item\hypertarget{radionbib:WheelerFeynman1945}{}
J.~A.~Wheeler y R.~P.~Feynman,
\emph{Interaction with the Absorber as the Mechanism of Radiation},
Reviews of Modern Physics 17 (1945), 157--181.
\url{https://doi.org/10.1103/RevModPhys.17.157}

\item\hypertarget{radionbib:Feynman1948}{}
R.~P.~Feynman,
\emph{Space-Time Approach to Non-Relativistic Quantum Mechanics},
Reviews of Modern Physics 20 (1948), 367--387.
\url{https://doi.org/10.1103/RevModPhys.20.367}

\item\hypertarget{radionbib:Feynman1949}{}
R.~P.~Feynman,
\emph{The Theory of Positrons},
Physical Review 76 (1949), 749--759.
\url{https://doi.org/10.1103/PhysRev.76.749}
\end{enumerate}

\begin{narrativebox}
Los fundadores contienen las piezas: acción, frecuencia angular, variables
conjugadas, oscilador, incertidumbre, fase de caminos y reversión temporal.
La aportación del radión es reunirlas en un objeto generado: una elipse de
área \(h\), un sector de área \(\hbar\), un interior de profundidad infinita,
una orientación de hoja y un acarreo exacto de compatibilidad.
\end{narrativebox}
% Fuente: ../incoming/radion_monografia_20260827/manuscrito/sections/E_conclusion.tex:1-107
\subsubsection{Síntesis final: unidad angular con genealogía de fase, acción y memoria}

El radián no desaparece. Se vuelve visible por completo.

En su uso ordinario, un radián es una razón que no necesita recordar cómo fue
producida. El arco y el radio se cancelan dimensionalmente; la unidad parece
desnuda. Esa desnudez es una virtud para la trigonometría, pero es un olvido
para una teoría que pretende explicar la estructura del continuo y el origen
de la acción. El radión recupera lo olvidado sin alterar el valor publicado.

La genealogía comienza en APP, donde una salida conserva valor y también
hoja, cociente, resto, acarreo y frontera. El TRIT introduce régimen y
orientación. El TPK hace volver la fase sin reiniciar el estado. La conexión
nonádica construye la sección coinductiva de \(\pi\), y la dodecafase publica
la bisagra de \(50^\circ\). El par angular aporta un modo común \(A\) y una
apertura orientada \(C^*\). La incidencia de seis fases activas produce
\(90\), sus dos hojas sobre el ambiente \(729\) producen \(20/81\), y la
torsión global aporta la segunda sección.

Entonces aparece \(\varepsilon_R\). No como un decimal suplicado por el
cierre, sino como la diferencia que la estructura estaba obligada a
transportar:
\[
\epsone=\frac{20}{81}\Delta_4-(S_E-1).
\]
El acarreo publica sus cifras; la profundidad certifica su límite; la carta
angular lo reconoce como
\[
5.13830167585752820716851646226459474899\ldots
\times10^{-8}\,{}^\circ.
\]
La ecuación del radián queda cerrada porque ambos lados son publicaciones del
mismo estado:
\BigRadionEquation

La historia no termina en la igualdad. El par \((A,C^*)\) construye una
elipse positiva. Su producto de semiejes fija \(2\hret\); su cociente fija la
forma. El teorema de área demuestra que una unidad de fase barre \(\hret\) y
que la vuelta encierra \(\hfull\). Aquí adquiere contenido físico la división
por \(2\pi\): \(h\) es acción de vuelta; \(\hbar\) es acción por radión.

Los focos dejan de ser adornos de una figura. Su separación reconstruye
\(C^*\) en la carta espectral, y su distancia de Klein mide una profundidad
hiperbólica finita dentro de un dominio cuya área hiperbólica total es
infinita. La misma frontera contiene una acción finita. La aparente paradoja
es la tesis central de la geometría: una medida cierra la acción; otra abre la
profundidad.

El círculo, por tanto, no es falso. Es la proyección exacta que conserva fase
y olvida anisotropía, interior y memoria. La elipse conserva acción y forma.
Klein conserva profundidad. La hélice conserva retorno y biografía. Ninguna
imagen sustituye a las demás; todas nacen del mismo estado.

La realización de Hilbert reconoce la elipse como contorno de covarianza
mínima. El oscilador LC la reconoce como intercambio de reservas eléctrica y
magnética. El lector \(90/120\) publica permeabilidad, permitividad,
impedancia y velocidad. La incidencia seis--ocho produce una pantalla causal;
Hodge convierte orientación en dualidad de campos; Lorentz convierte campo y
carga en dinámica; Cartan--Holst recibe la holonomía en geometría
gravitatoria. El orden importa: el objeto HMT selecciona; la física
convencional realiza.

La palabra \emph{radión} queda así justificada. No significa que un foco sea
una carga positiva y el otro una carga negativa. Significa que la unidad de
fase posee orientación firmada antes de la publicación eléctrica:
\[
\mathcal A(\mathfrak r_+)=+\hret,
\qquad
\mathcal A(\mathfrak r_-)=-\hret.
\]
La carga aparece cuando Mecánica Dimensional aplica el transductor. El signo
no se añade al final; se conserva desde la hoja.

También queda explicado por qué el sector \(90/120\) parecía esconder muchas
respuestas. Incidencia, extractor dodecafásico, cola espectral, núcleo de
Barbero, torres y acarreo comparten origen. Precisamente por compartir origen
pueden compararse. Precisamente por tener operaciones distintas no deben
fundirse. La tipificación no reduce la unidad del sistema: la hace inteligible.

La conclusión más compacta de esta obra puede escribirse en cuatro líneas:
\[
\boxed{
\begin{aligned}
\text{fase:}\quad&\theta=1,\\
\text{acción:}\quad&\mathcal A=\sigma\hret,\\
\text{memoria:}\quad&H_9(n+9)=H_9(n)+(0,1),\\
\text{compatibilidad:}\quad&1=S_E-\Phi_T+\epsone.
\end{aligned}}
\]

El radión es esa unidad cuádruple. El radián es su sombra exacta en la carta
angular. La diferencia no está en el valor de \(\pi\); está en cuánta
estructura estamos dispuestos a conservar cuando decimos que una vuelta ha
cerrado.

\vspace{2\baselineskip}
\begin{center}
\begin{minipage}{.82\textwidth}
\centering\large\itshape\color{RadNavy}
El círculo publica la fase.\\
La elipse contiene la acción.\\
Klein abre la profundidad.\\
La hélice recuerda el camino.\\[.7em]
El radión es la unidad que no obliga a elegir entre ellas.
\end{minipage}
\end{center}

